\section{Homogeneous prime spectra}
\label{section:II.2}

\subsection{Generalities on graded rings and modules.}
\label{subsection:II.2.1}

\begin{notation}[2.1.1]
\label{II.2.1.1}
Given a ring $S$ \emph{graded} in positive degrees, we denote
by $S_n$ the subset of $S$ consisting of the homogeneous elements of degree
$n$ ($n\geq0$), and by $S_+$ the (direct) sum of the $S_n$ such that $n>0$;
we have $1\in S_0$, $S_0$ is a subring of $S$, $S_+$ is a graded ideal of
$S$, and $S$ is the direct sum of $S_0$ and $S_+$. If $M$ is a \emph{graded
module} over $S$ (with positive or negative degrees), we similarly denote by
$M_n$ the $S_0$-module consisting of the homogeneous elements of degree $n$
of $M$ (with $n\in\mathbf{Z}$).

For every integer $d>0$, we denote by $S^{(d)}$ the direct sum of the
$S_{nd}$; regarding the elements of $S_{nd}$ as homogeneous of degree
$n$, the $S_{nd}$ define a graded ring structure on $S^{(d)}$.

For every integer $k$ such that $0\leq k\leq d-1$, we denote by
$M^{(d,k)}$ the direct sum
\oldpage[II]{20}
of the $M_{nd+k}$ ($n\in\mathbf{Z}$); it is a graded $S^{(d)}$-module when
the elements of $M_{nd+k}$ are regarded as homogeneous of degree $n$. We
write $M^{(d)}$ instead of $M^{(d,0)}$.

With the preceding notation, for every integer $n$ (positive or negative),
we denote by $M(n)$ the graded $S$-module defined by
$(M(n))_k=M_{n+k}$ for every $k\in\mathbf{Z}$. In particular, $S(n)$ will be
a graded $S$-module such that $(S(n))_k=S_{n+k}$, with the convention
$S_n=0$ for $n<0$. We say that a graded $S$-module $M$ is \emph{free}
if it is isomorphic, as a \emph{graded module}, to a direct sum of
modules of the form $S(n)$; since $S(n)$ is a monogenic $S$-module,
generated by the element $1$ of $S$ regarded as an element of degree $-n$,
it amounts to the same to say that $M$ admits a \emph{basis} over $S$
consisting of \emph{homogeneous} elements.

We say that a graded $S$-module $M$ \emph{admits a finite presentation} if
there exists an exact sequence $P\to Q\to M\to0$, where $P$ and $Q$ are
finite direct sums of modules of the form $S(n)$ and the homomorphisms are
of degree $0$ (cf. (\hyperref[II.2.1.2]{2.1.2})).
\end{notation}

\begin{env}[2.1.2]
\label{II.2.1.2}
Let $M$, $N$ be two graded $S$-modules; we define on $M\otimes_S N$ a
structure of \emph{graded} $S$-module as follows. On the tensor product
$M\otimes_{\mathbf{Z}}N$, we can define a structure of graded
$\mathbf{Z}$-module (where $\mathbf{Z}$ is graded by
$\mathbf{Z}_0=\mathbf{Z}$, $\mathbf{Z}_n=0$ for $n\ne0$) by setting
\[
  (M\otimes_{\mathbf{Z}}N)_q
  =\bigoplus_{m+n=q}M_m\otimes_{\mathbf{Z}}N_n
\]
(since $M$ and $N$ are the direct sums of the $M_m$ and of the $N_n$
respectively, we know that $M\otimes_{\mathbf{Z}}N$ can be canonically
identified with the direct sum of all the
$M_m\otimes_{\mathbf{Z}}N_n$). This being so, we have
$M\otimes_S N=(M\otimes_{\mathbf{Z}}N)/P$, where $P$ is the
sub-$\mathbf{Z}$-module of $M\otimes_{\mathbf{Z}}N$ generated by the
elements $(xs)\otimes y-x\otimes(sy)$ for $x\in M$, $y\in N$, $s\in S$;
it is clear that $P$ is a \emph{graded} sub-$\mathbf{Z}$-module of
$M\otimes_{\mathbf{Z}}N$, and we see immediately that passing to the
quotient indeed gives a structure of graded $S$-module on $M\otimes_S N$.

For two graded $S$-modules $M$, $N$, recall that a homomorphism
$u:M\to N$ of $S$-modules is said to be \emph{of degree $k$} if
$u(M_j)\subset N_{j+k}$ for every $j\in\mathbf{Z}$. If $H_n$ denotes the
set of all homomorphisms of degree $n$ from $M$ to $N$, we denote by
$\operatorname{Hom}_S(M,N)$ the (direct) \emph{sum} of the $H_n$
($n\in\mathbf{Z}$) in the $S$-module $H$ of all homomorphisms (of
$S$-modules) from $M$ to $N$; in general $\operatorname{Hom}_S(M,N)$ is not
equal to the latter. However, we have $H=\operatorname{Hom}_S(M,N)$ when
$M$ is \emph{of finite type}; indeed, we may then suppose $M$ generated by
finitely many homogeneous elements $x_i$ ($1\leq i\leq n$), and every
homomorphism $u\in H$ can be written in exactly one way as
$\sum_{k\in\mathbf{Z}}u_k$, where for each $k$, $u_k(x_i)$ is equal to the
homogeneous component of degree $k+\deg(x_i)$ of $u(x_i)$ ($1\leq i\leq n$),
which implies $u_k=0$ except for finitely many indices; by definition we
have $u_k\in H_k$, whence the conclusion.

We say that the elements of degree $0$ of $\operatorname{Hom}_S(M,N)$ are
\emph{homomorphisms of graded $S$-modules}. It is clear that
$S_mH_n\subset H_{m+n}$, so the $H_n$ define a structure of graded
$S$-module on $\operatorname{Hom}_S(M,N)$.

It follows immediately from these definitions that we have
\[
  \label{II.2.1.2.1}
  M(m)\otimes_S N(n)=(M\otimes_S N)(m+n),
  \tag{2.1.2.1}
\]
\[
  \label{II.2.1.2.2}
  \operatorname{Hom}_S(M(m),N(n))
  =(\operatorname{Hom}_S(M,N))(n-m)
  \tag{2.1.2.2}
\]
for two graded $S$-modules $M$, $N$.
\oldpage[II]{21}
Let $S$, $S'$ be two graded rings; a homomorphism of \emph{graded
rings} $\varphi:S\to S'$ is a ring homomorphism such that
$\varphi(S_n)\subset S'_n$ for every $n\in\mathbf{Z}$ (in other words,
$\varphi$ must be a homomorphism \emph{of degree $0$} of graded
$\mathbf{Z}$-modules). Such a homomorphism defines on $S'$ a structure of
\emph{graded} $S$-module; endowed with this structure and with its graded
ring structure, $S'$ is said to be a \emph{graded $S$-algebra}.

If then $M$ is a graded $S$-module, the tensor product $M\otimes_S S'$
of graded $S$-modules is naturally endowed with a structure of
\emph{graded} $S'$-module, the grading being the one defined above.
\end{env}

\begin{lemma}[2.1.3]
\label{II.2.1.3}
Let $S$ be a ring graded in positive degrees. For a subset $E$ of $S_+$
consisting of homogeneous elements to generate $S_+$ as an $S$-module, it is
necessary and sufficient that $E$ generate $S$ as an $S_0$-algebra.
\end{lemma}

\begin{proof}
The condition is evidently sufficient; let us show that it is necessary.
Let $E_n$ (resp. $E^n$) be the set of elements of $E$ of degree equal to $n$
(resp. $\leq n$); it will suffice to show, by induction on $n>0$, that $S_n$
is the $S_0$-module generated by the elements of degree $n$ which are
products of elements of $E^n$. Now this is evident for $n=1$ by virtue of
the hypothesis; the latter moreover shows that
$S_n=\sum_{p=0}^{n-1}S_pE_{n-p}$, and the induction argument is then
immediate.
\end{proof}

\begin{corollary}[2.1.4]
\label{II.2.1.4}
For $S_+$ to be an ideal of finite type, it is necessary and sufficient that
$S$ be an $S_0$-algebra of finite type.
\end{corollary}

\begin{proof}
Indeed, we can always suppose that a finite system of generators of the
$S_0$-algebra $S$ (resp. of the $S$-ideal $S_+$) consists of homogeneous
elements, by replacing each of the generators considered by its homogeneous
components.
\end{proof}

\begin{corollary}[2.1.5]
\label{II.2.1.5}
For $S$ to be Noetherian, it is necessary and sufficient that $S_0$ be
Noetherian and that $S$ be an $S_0$-algebra of finite type.
\end{corollary}

\begin{proof}
The condition is evidently sufficient; it is necessary, since $S_0$
is isomorphic to $S/S_+$ and $S_+$ must be an ideal of finite type
(\hyperref[II.2.1.4]{2.1.4}).
\end{proof}

\begin{lemma}[2.1.6]
\label{II.2.1.6}
Let $S$ be a ring graded in positive degrees which is an $S_0$-algebra of
finite type. Let $M$ be a graded $S$-module of finite type. Then:
\begin{enumerate}
  \item[\rm(i)] The $M_n$ are $S_0$-modules of finite type, and there exists
    an integer $n_0$ such that $M_n=0$ for $n\leq n_0$.
  \item[\rm(ii)] There exist an integer $n_1$ and an integer $h>0$ such that,
    for every integer $n\geq n_1$, we have $M_{n+h}=S_hM_n$.
  \item[\rm(iii)] For every pair of integers $(d,k)$ such that $d>0$,
    $0\leq k\leq d-1$, $M^{(d,k)}$ is an $S^{(d)}$-module of finite type.
  \item[\rm(iv)] For every integer $d>0$, $S^{(d)}$ is an $S_0$-algebra of
    finite type.
  \item[\rm(v)] There exists an integer $h>0$ such that $S_{mh}=(S_h)^m$ for
    every $m>0$.
  \item[\rm(vi)] For every integer $n>0$, there exists an integer $m_0$ such
    that $S_m\subset S_+^n$ for every $m\geq m_0$.
\end{enumerate}
\end{lemma}

\begin{proof}
We may suppose $S$ generated (as an $S_0$-algebra) by homogeneous
elements $f_i$ of degree $h_i$ ($1\leq i\leq r$), and $M$ generated (as
an $S$-module) by homogeneous elements $x_j$ of degree $k_j$
($1\leq j\leq s$). It is clear that $M_n$ consists of the linear
\oldpage[II]{22}
combinations, with coefficients in $S_0$, of the elements
$f_1^{\alpha_1}\cdots f_r^{\alpha_r}x_j$ such that the $\alpha_i$ are
integers $\geq 0$ satisfying $k_j+\sum_i\alpha_i h_i=n$; for each $j$,
there are only finitely many systems $(\alpha_i)$ satisfying this equation,
since the $h_i$ are $>0$, whence the first assertion of (i); the second
is evident. On the other hand, let $h$ be the l.c.m. of the $h_i$ and set
$g_i=f_i^{h/h_i}$ ($1\leq i\leq r$), so that all the $g_i$ are of
degree $h$; let $z_\mu$ be the elements of $M$ of the form
$f_1^{\alpha_1}\cdots f_r^{\alpha_r}x_j$ with
$0\leq\alpha_i<h/h_i$ for $1\leq i\leq r$; these elements are finite in
number; let $n_1$ be the largest of their degrees. It is clear that for
$n\geq n_1$, every element of $M_{n+h}$ is a linear combination of the
$z_\mu$ whose coefficients are monomials of degree $>0$ in the $g_i$,
hence we have $M_{n+h}=S_hM_n$, which establishes (ii). In the same way, we
see (for every $d>0$) that an element of $M^{(d,k)}$ is a linear
combination, with coefficients in $S_0$, of elements of the form
$g^df_1^{\alpha_1}\cdots f_r^{\alpha_r}x_j$ with
$0\leq\alpha_i<d$, $g$ being a homogeneous element of $S$; whence (iii); (iv)
then follows from (iii) and from Lemma (\hyperref[II.2.1.3]{2.1.3}),
taking $M=S_+$, since $(S_+)^{(d)}=(S^{(d)})_+$. Assertion (v)
follows from (ii) by taking $M=S$. Finally, for a given $n$, there are only
finitely many systems $(\alpha_i)$ such that $\alpha_i\geq 0$ and
$\sum_i\alpha_i<n$, so if $m_0$ is the largest value of the sum
$\sum_i\alpha_i h_i$ for these systems, we have $S_m\subset S_+^n$ for
$m>m_0$, which proves (vi).
\end{proof}

\begin{corollary}[2.1.7]
\label{II.2.1.7}
If $S$ is Noetherian, so is $S^{(d)}$ for every integer
$d>0$.
\end{corollary}

\begin{proof}
This follows from (\hyperref[II.2.1.5]{2.1.5}) and
(\hyperref[II.2.1.6]{2.1.6, (iv)}).
\end{proof}

\begin{env}[2.1.8]
\label{II.2.1.8}
Let $\mathfrak{p}$ be a graded prime ideal of the graded ring $S$;
$\mathfrak{p}$ is therefore the direct sum of the subgroups
$\mathfrak{p}_n=\mathfrak{p}\cap S_n$. Suppose that $\mathfrak{p}$ does not
contain $S_+$. Then, if $f\in S_+$ does not belong to
$\mathfrak{p}$, the relation $f^nx\in\mathfrak{p}$ is equivalent to
$x\in\mathfrak{p}$; in particular, if $f\in S_d$ ($d>0$), for every
$x\in S_{m-nd}$, the relation $f^nx\in\mathfrak{p}_m$ is equivalent to
$x\in\mathfrak{p}_{m-nd}$.
\end{env}

\begin{proposition}[2.1.9]
\label{II.2.1.9}
Let $n_0$ be an integer $>0$; for every $n\geq n_0$ let $\mathfrak{p}_n$ be
a subgroup of $S_n$. For there to exist a graded prime ideal
$\mathfrak{p}$ of $S$ not containing $S_+$ and such that
$\mathfrak{p}\cap S_n=\mathfrak{p}_n$ for every $n\geq n_0$, it is necessary
and sufficient that the following conditions be satisfied:
\begin{enumerate}
  \item[$1^\circ$] $S_m\mathfrak{p}_n\subset\mathfrak{p}_{m+n}$ for every
    $m\geq 0$ and every $n\geq n_0$.
  \item[$2^\circ$] For $m\geq n_0$, $n\geq n_0$, $f\in S_m$, $g\in S_n$,
    the relation $fg\in\mathfrak{p}_{m+n}$ implies
    $f\in\mathfrak{p}_m$ or $g\in\mathfrak{p}_n$.
  \item[$3^\circ$] $\mathfrak{p}_n\neq S_n$ for at least one $n\geq n_0$.
\end{enumerate}
Moreover, the graded prime ideal $\mathfrak{p}$ is then unique.
\end{proposition}

\begin{proof}
It is evident that conditions $1^\circ$ and $2^\circ$ are necessary.
Moreover, if $\mathfrak{p}\not\supset S_+$, there exists at least one $k>0$
such that $\mathfrak{p}\cap S_k\neq S_k$; if $f\in S_k$ does not belong to
$\mathfrak{p}$, the relation $\mathfrak{p}\cap S_n=S_n$ implies
$\mathfrak{p}\cap S_{n-mk}=S_{n-mk}$ by
(\hyperref[II.2.1.8]{2.1.8}); hence, if
$\mathfrak{p}\cap S_n=S_n$ from some value of $n$ onwards, we would have
$\mathfrak{p}\supset S_+$, contrary to the hypothesis, which proves that
$3^\circ$ is necessary. Conversely, suppose conditions
$1^\circ$, $2^\circ$ and $3^\circ$ are satisfied. Note that if, for an
integer $d\geq n_0$, $f\in S_d$ does not belong to $\mathfrak{p}_d$, then,
if $\mathfrak{p}$ exists, $\mathfrak{p}_m$, for $m<n_0$, is necessarily
equal to the set of $x\in S_m$ such that
$f^rx\in\mathfrak{p}_{m+rd}$ except for finitely many values of $r$.
This already proves that if $\mathfrak{p}$ exists, it is unique. It remains
to show that if we define the $\mathfrak{p}_m$ for $m<n_0$ by the preceding
condition, then $\mathfrak{p}=\sum_{n=0}^{\infty}\mathfrak{p}_n$ is a prime
ideal. Note first that by virtue of $2^\circ$, for $m\geq n_0$,
$\mathfrak{p}_m$ is also defined as the set of $x\in S_m$ such that
$f^rx\in\mathfrak{p}_{m+rd}$ except for finitely many values of $r$.
This
\oldpage[II]{23}
being so, if $g\in S_m$, $x\in\mathfrak{p}_n$, we have
$f^rgx\in\mathfrak{p}_{m+n+rd}$ except for finitely many values of $r$,
hence $gx\in\mathfrak{p}_{m+n}$, which proves that $\mathfrak{p}$ is an
ideal of $S$. To establish that this ideal is prime, in other words that
the ring $S/\mathfrak{p}$, graded by the subgroups
$S_n/\mathfrak{p}_n$, is an integral ring, it suffices (by considering the
components of highest degree of two elements of $S/\mathfrak{p}$) to
prove that if $x\in S_m$, $y\in S_n$ are such that
$x\notin\mathfrak{p}_m$, $y\notin\mathfrak{p}_n$, then
$xy\notin\mathfrak{p}_{m+n}$. Otherwise, for $r$ sufficiently large, we would
have $f^{2r}xy\in\mathfrak{p}_{m+n+2rd}$; but we have
$f^ry\notin\mathfrak{p}_{n+rd}$ for every $r>0$; it then follows from
$2^\circ$ that, except for finitely many values of $r$, we have
$f^rx\in\mathfrak{p}_{m+rd}$, and we would conclude that
$x\in\mathfrak{p}_m$, contrary to the hypothesis.
\end{proof}

\begin{env}[2.1.10]
\label{II.2.1.10}
We will say that a subset $\mathfrak{J}$ of $S_+$ is an
\emph{ideal of $S_+$} if it is an ideal of $S$, and that $\mathfrak{J}$ is a
\emph{graded prime ideal of $S_+$} if it is the intersection of $S_+$ and
a graded prime ideal of $S$ \emph{not containing $S_+$} (this prime ideal
is moreover unique by
(\hyperref[II.2.1.9]{2.1.9})). If $\mathfrak{J}$ is an ideal of $S_+$,
the \emph{radical of $\mathfrak{J}$ in $S_+$} is the set of elements of
$S_+$ some power of which belongs to $\mathfrak{J}$; in other words, it is
the set $r_+(\mathfrak{J})=r(\mathfrak{J})\cap S_+$; in particular, the
radical of $0$ in $S_+$ is also called the \emph{nilradical} of $S_+$ and
denoted by $\mathfrak{N}_+$: it is thus the set of nilpotent elements of
$S_+$. If $\mathfrak{J}$ is a \emph{graded} ideal of $S_+$, its radical
$r_+(\mathfrak{J})$ is a graded ideal: passing to the quotient ring
$S/\mathfrak{J}$, we can reduce to the case $\mathfrak{J}=0$, and
everything comes down to seeing that if $x=x_h+x_{h+1}+\cdots+x_k$ is
nilpotent, then so are the $x_i\in S_i$ ($1\leq h\leq i\leq k$); we may
suppose $x_k\neq 0$, and the component of largest degree of $x^n$ is then
$x_k^n$, so $x_k$ is nilpotent, and we then argue by induction on $k$. We
say that the graded ring $S$ is \emph{essentially reduced} if
$\mathfrak{N}_+=0$, in other words if $S_+$ contains no nilpotent elements
$\neq 0$.
\end{env}

\begin{env}[2.1.11]
\label{II.2.1.11}
Note that if, in the graded ring $S$, an element $x$ is a divisor of
$0$, then so is its component of highest degree. We say that a ring $S$ is
\emph{essentially integral} if the ring $S_+$
(\emph{without unit element}) contains no divisor of $0$ and is
$\neq 0$; for this it suffices that no homogeneous element $\neq 0$ in
$S_+$ be a divisor of $0$ in this ring. It is clear that if $\mathfrak{p}$
is a graded prime ideal of $S_+$, then $S/\mathfrak{p}$ is essentially
integral.

Let $S$ be an essentially integral graded ring, and let $x_0\in S_0$;
if there then exists \emph{one} homogeneous element $f\neq 0$ of $S_+$ such
that $x_0f=0$, we have $x_0S_+=0$, for we have $(x_0g)f=(x_0f)g=0$ for
every $g\in S_+$, and the hypothesis therefore implies $x_0g=0$. Hence,
for $S$ to be integral, it is necessary and sufficient that $S_0$ be
integral and that the annihilator of $S_+$ in $S_0$ consist of $0$ alone.
\end{env}

\subsection{Rings of fractions of a graded ring.}
\label{subsection:II.2.2}

\begin{env}[2.2.1]
\label{II.2.2.1}
Let $S$ be a ring graded in positive degrees, and $f$ a
\emph{homogeneous} element of $S$ of degree $d>0$; then the ring of
fractions $S'=S_f$ is graded, taking for $S'_n$ the set of $x/f^k$, where
$x\in S_{n+kd}$ with $k\geq 0$ (observe here that $n$ can take arbitrary
negative values); we denote the subring $S'_0=(S_f)_0$ of $S'$
consisting of the elements \emph{of degree $0$} by
$S_{(f)}$.

If $f\in S_d$, the monomials $(f/1)^h$ in $S_f$ ($h$ a positive or
negative integer) form a \emph{linearly independent family} over the ring
$S_{(f)}$, and the set of their linear combinations is none other than
\oldpage[II]{24}
the ring $(S^{(d)})_f$, which is therefore \emph{isomorphic to
$S_{(f)}[T,T^{-1}]=S_{(f)}\otimes_{\mathbf{Z}}\mathbf{Z}[T,T^{-1}]$}
(where $T$ is an indeterminate). Indeed, if we have a relation
$\sum_{h=-a}^{b}z_h(f/1)^h=0$ with $z_h=x_h/f^m$, where all the $x_h$
belong to $S_{md}$, this relation is equivalent by definition to
the existence of a $k\geq a$ such that
$\sum_{h=-a}^{b}f^{h+k}x_h=0$, and since the degrees of the terms of this
sum are distinct, we have $f^{h+k}x_h=0$ for every $h$, whence $z_h=0$ for
every $h$.

If $M$ is a graded $S$-module, $M'=M_f$ is a graded $S_f$-module,
$M'_n$ being the set of $z/f^k$ with $z\in M_{n+kd}$ ($k\geq 0$); we
denote by $M_{(f)}$ the set of homogeneous elements of degree $0$ of
$M'$; it is immediate that $M_{(f)}$ is an $S_{(f)}$-module and that we have
$(M^{(d)})_f=M_{(f)}\otimes_{S_{(f)}}(S^{(d)})_f$.
\end{env}

\begin{lemma}[2.2.2]
\label{II.2.2.2}
Let $d$, $e$ be two integers $>0$, $f\in S_d$, $g\in S_e$. There exists a
canonical ring isomorphism
\[
  S_{(fg)}\xrightarrow{\sim}(S_{(f)})_{g^d/f^e}~;
\]
if we canonically identify these two rings, there exists a canonical module
isomorphism
\[
  M_{(fg)}\xrightarrow{\sim}(M_{(f)})_{g^d/f^e}.
\]
\end{lemma}

\begin{proof}
Indeed, $fg$ divides $f^e g^d$, and the latter element divides $(fg)^{de}$,
so the graded rings $S_{fg}$ and $S_{f^e g^d}$ are canonically
identified; on the other hand, $S_{f^e g^d}$ is also identified with
$(S_{f^e})_{g^d/1}$ (\hyperref[0.1.4.6]{0, I.4.6}), and since $f^e/1$ is
invertible in $S_{f^e}$, $S_{f^e g^d}$ is also identified with
$(S_{f^e})_{g^d/f^e}$. Now the element $g^d/f^e$ is of degree $0$ in
$S_{f^e}$; we conclude immediately that the subring of
$(S_{f^e})_{g^d/f^e}$ consisting of the elements of degree $0$ is
$(S_{(f^e)})_{g^d/f^e}$, and since evidently
$S_{(f^e)}=S_{(f)}$, this proves the first part of the proposition;
the second is established in the same way.
\end{proof}

\begin{env}[2.2.3]
\label{II.2.2.3}
With the hypotheses of (\hyperref[II.2.2.2]{2.2.2}), it is clear that
the canonical homomorphism $S_f\to S_{fg}$
(\hyperref[0.1.4.1]{0, I.4.1}), which maps $x/f^k$ to
$g^k x/(fg)^k$, is of degree $0$, and therefore gives by restriction a
\emph{canonical homomorphism $S_{(f)}\to S_{(fg)}$} such that the diagram
\[
  \xymatrix{
    & S_{(f)}\ar[dl]\ar[dr]\\
    S_{(fg)}\ar[rr]^-{\sim} && (S_{(f)})_{g^d/f^e}
  }
\]
is commutative. We define in the same way a canonical homomorphism
$M_{(f)}\to M_{(fg)}$.
\end{env}

\begin{lemma}[2.2.4]
\label{II.2.2.4}
If $f$, $g$ are two homogeneous elements of $S_+$, the ring $S_{(fg)}$ is
generated by the union of the canonical images of $S_{(f)}$ and $S_{(g)}$.
\end{lemma}

\begin{proof}
By virtue of (\hyperref[II.2.2.2]{2.2.2}), it suffices to see that
$1/(g^d/f^e)=f^{d+e}/(fg)^d$ belongs to the canonical image of $S_{(g)}$
in $S_{(fg)}$, which is evident by definition.
\end{proof}

\begin{proposition}[2.2.5]
\label{II.2.2.5}
Let $d$ be an integer $>0$ and let $f\in S_d$. Then there exists a canonical
ring isomorphism
$S_{(f)}\xrightarrow{\sim}S^{(d)}/(f-1)S^{(d)}$; if we identify these two
rings by this isomorphism, there exists a canonical module isomorphism
$M_{(f)}\xrightarrow{\sim}M^{(d)}/(f-1)M^{(d)}$.
\end{proposition}

\begin{proof}
The first of these isomorphisms is defined by assigning to
$x/f^n$, where $x\in S_{nd}$, the element $\overline{x}$, the class of $x$
mod. $(f-1)S^{(d)}$; this map is well defined, since we have the congruence
$f^h x\equiv x\pmod{(f-1)S^{(d)}}$ for every $x\in S^{(d)}$, so if
$f^h x=0$ for some $h>0$,
\oldpage[II]{25}
we have $\overline{x}=0$. On the other hand, if $x\in S_{nd}$ is such that
$x=(f-1)y$ with
$y=y_{hd}+y_{(h+1)d}+\cdots+y_{kd}$, where $y_{jd}\in S_{jd}$ and
$y_{hd}\neq 0$, we necessarily have $h=n$ and $x=-y_{hd}$, as well as the
relations $y_{(j+1)d}=fy_{jd}$ for $h\leq j\leq k-1$, $fy_{kd}=0$, which
finally gives $f^{k-n+1}x=0$; to every class $\overline{x}$ mod.
$(f-1)S^{(d)}$ of an element $x\in S_{nd}$ we can therefore assign
the element $x/f^n$ of $S_{(f)}$, since the preceding remark shows that this
map is well defined. It is immediate that the two maps thus
defined are ring homomorphisms, inverse to each other. We
proceed in exactly the same way for $M$.
\end{proof}

\begin{corollary}[2.2.6]
\label{II.2.2.6}
If $S$ is Noetherian, so is $S_{(f)}$ for every homogeneous $f$
of degree $>0$.
\end{corollary}

\begin{proof}
This follows immediately from (\hyperref[II.2.1.7]{2.1.7}) and
(\hyperref[II.2.2.5]{2.2.5}).
\end{proof}

\begin{env}[2.2.7]
\label{II.2.2.7}
Let $T$ be a multiplicative subset of $S_+$ consisting of
\emph{homogeneous} elements; $T_0=T\cup\{1\}$ is then a multiplicative
subset of $S$; since the elements of $T_0$ are homogeneous, the ring
$T_0^{-1}S$ is again graded in an evident way; we denote by $S_{(T)}$ the
subring of $T_0^{-1}S$ consisting of the elements of order $0$, that is,
the elements of the form $x/h$, where $h\in T$ and $x$ is homogeneous of
degree equal to that of $h$. We know (\hyperref[0.1.4.5]{0, 1.4.5}) that
$T_0^{-1}S$ is canonically identified with the inductive limit of the rings
$S_f$, where $f$ runs through $T$ (for the canonical homomorphisms
$S_f\to S_{fg}$); since this identification respects degrees, it identifies
$S_{(T)}$ with the \emph{inductive limit} of the $S_{(f)}$ for $f\in T$.
For every graded $S$-module $M$, we define in the same way the module
$M_{(T)}$ (over the ring $S_{(T)}$) consisting of the elements of degree $0$
of $T_0^{-1}M$, and we see that this module is the inductive limit of the
$M_{(f)}$ for $f\in T$.

If $\mathfrak{p}$ is a graded prime ideal of $S_+$, we denote by
$S_{(\mathfrak{p})}$ and $M_{(\mathfrak{p})}$ the ring $S_{(T)}$ and the
module $M_{(T)}$ respectively, where $T$ is the set of \emph{homogeneous}
elements of $S_+$ which do not belong to $\mathfrak{p}$.
\end{env}

\subsection{Homogeneous prime spectrum of a graded ring.}
\label{subsection:II.2.3}

\begin{env}[2.3.1]
\label{II.2.3.1}
Given a ring $S$ graded in positive degrees, we call
\emph{homogeneous prime spectrum} of $S$, and denote by
$\operatorname{Proj}(S)$, the set of graded prime ideals of $S_+$
(\hyperref[II.2.1.10]{2.1.10}), or, what amounts to the same, the set
of graded prime ideals of $S$ \emph{not containing} $S_+$; we
are going to define a \emph{scheme} structure having
$\operatorname{Proj}(S)$ as underlying set.
\end{env}

\begin{env}[2.3.2]
\label{II.2.3.2}
For every subset $E$ of $S$, let $V_+(E)$ be the set of graded prime
ideals of $S$ containing $E$ and not containing $S_+$; it is thus the
subset $V(E)\cap\operatorname{Proj}(S)$ of
$\operatorname{Spec}(S)$. From (\hyperref[I.1.1.2]{I, 1.1.2}) we
therefore deduce:
\[
  \label{II.2.3.2.1}
  V_+(0)=\operatorname{Proj}(S),\quad
  V_+(S)=V_+(S_+)=\varnothing
  \tag{2.3.2.1}
\]
\[
  \label{II.2.3.2.2}
  V_+\left(\bigcup_\lambda E_\lambda\right)
  =\bigcap_\lambda V_+(E_\lambda)
  \tag{2.3.2.2}
\]
\[
  \label{II.2.3.2.3}
  V_+(EE')=V_+(E)\cup V_+(E')
  \tag{2.3.2.3}
\]
$V_+(E)$ does not change when $E$ is replaced by the graded ideal generated
by $E$; moreover, if $\mathfrak{J}$ is a graded ideal of $S$, we have
\[
  \label{II.2.3.2.4}
  V_+(\mathfrak{J})
  =V_+\left(\bigcup_{q\geq n}(\mathfrak{J}\cap S_q)\right)
  \tag{2.3.2.4}
\]
\oldpage[II]{26}
for every $n>0$: indeed, if
$\mathfrak{p}\in\operatorname{Proj}(S)$ contains the homogeneous elements of
$\mathfrak{J}$ of degree $\geq n$, then, since by hypothesis there exists a
homogeneous element $f\in S_d$ not contained in $\mathfrak{p}$, for every
$m\geq 0$ and every $x\in S_m\cap\mathfrak{J}$ we have
$f^r x\in\mathfrak{J}\cap S_{m+rd}$ for every $r\geq0$; except for
finitely many values of $r$, we have $m+rd\geq n$, hence
$f^r x\in\mathfrak{p}\cap S_{m+rd}$, which implies
$x\in\mathfrak{p}\cap S_m$
(\hyperref[II.2.1.8]{2.1.8}).

Finally, for every graded ideal $\mathfrak{J}$ of $S$, we have
\[
  \label{II.2.3.2.5}
  V_+(\mathfrak{J})=V_+(r_+(\mathfrak{J})).
  \tag{2.3.2.5}
\]
\end{env}

\begin{env}[2.3.3]
\label{II.2.3.3}
By definition, the $V_+(E)$ are the closed subsets of
$X=\operatorname{Proj}(S)$ for the topology induced by the spectral
topology of $\operatorname{Spec}(S)$, which we will also call the
\emph{spectral topology} on $X$. For every $f\in S$ we set
\[
  \label{II.2.3.3.1}
  D_+(f)=D(f)\cap\operatorname{Proj}(S)
  =\operatorname{Proj}(S)\setminus V_+(f)
  \tag{2.3.3.1}
\]
and consequently, for two elements $f$, $g$ of $S$
(\hyperref[I.1.1.9.1]{I, 1.1.9.1}), we have
\[
  \label{II.2.3.3.2}
  D_+(fg)=D_+(f)\cap D_+(g).
  \tag{2.3.3.2}
\]
\end{env}

\begin{proposition}[2.3.4]
\label{II.2.3.4}
When $f$ runs through the set of homogeneous elements of $S_+$, the
$D_+(f)$ form a basis of the topology of
$X=\operatorname{Proj}(S)$.
\end{proposition}

\begin{proof}
Indeed, it follows from (\hyperref[II.2.3.2.2]{2.3.2.2}) and
(\hyperref[II.2.3.2.4]{2.3.2.4}) that every closed subset of $X$ is an
intersection of sets of the form $V_+(f)$, where $f$ is homogeneous of
degree $>0$.
\end{proof}

\begin{env}[2.3.5]
\label{II.2.3.5}
Let $f$ be a \emph{homogeneous} element of $S_+$, of degree $d>0$; for
every graded prime ideal $\mathfrak{p}$ of $S$ not containing $f$, we know
that the set of $x/f^n$, where $x\in\mathfrak{p}$ and $n\geq 0$, is a
prime ideal of the ring of fractions $S_f$
(\hyperref[0.1.2.6]{0, 1.2.6}); its trace on $S_{(f)}$ is therefore a
prime ideal of this ring, which we denote by
$\psi_f(\mathfrak{p})$: it is the set of $x/f^n$ for $n\geq 0$,
$x\in\mathfrak{p}\cap S_{nd}$. We have thus defined a map
\[
  \psi_f:D_+(f)\longrightarrow\operatorname{Spec}(S_{(f)})~;
\]
moreover, if $g\in S_e$ is a second homogeneous element of $S_+$, we have a
commutative diagram
\[
  \label{II.2.3.5.1}
  \xymatrix{
    D_+(f)\ar[r]^{\psi_f}
    & \operatorname{Spec}(S_{(f)})\\
    D_+(fg)\ar[u]\ar[r]_{\psi_{fg}}
    & \operatorname{Spec}(S_{(fg)})\ar[u]
  }
  \tag{2.3.5.1}
\]
where the left vertical arrow is the inclusion, and the right one is
the map ${}^a\omega_{fg,f}$ deduced from the canonical homomorphism
$\omega=\omega_{fg,f}:S_{(f)}\to S_{(fg)}$
(\hyperref[I.1.2.1]{I, 1.2.1}). Indeed, if
$x/f^n\in\omega^{-1}(\psi_{fg}(\mathfrak{p}))$, where
$fg\notin\mathfrak{p}$, we have by definition
$g^n x/(fg)^n\in\psi_{fg}(\mathfrak{p})$, hence
$g^n x\in\mathfrak{p}$ and consequently $x\in\mathfrak{p}$, and the converse
is evident.
\end{env}

\begin{proposition}[2.3.6]
\label{II.2.3.6}
The map $\psi_f$ is a homeomorphism of $D_+(f)$ onto
$\operatorname{Spec}(S_{(f)})$.
\end{proposition}

\begin{proof}
In the first place, $\psi_f$ is continuous; for if $h\in S_{nd}$ is such
that $h/f^n\in\psi_f(\mathfrak{p})$, we have by definition
$h\in\mathfrak{p}$, and conversely, hence
\[
  \psi_f^{-1}(D(h/f^n))=D_+(hf)
\]
and our assertion follows from formula
(\hyperref[II.2.3.3.2]{2.3.3.2}). Moreover, the $D_+(hf)$, where $h$
runs through the sets $S_{nd}$, form a basis of the topology of
$D_+(f)$, by virtue of (\hyperref[II.2.3.4]{2.3.4}) and of formula
(\hyperref[II.2.3.3.2]{2.3.3.2}); the
\oldpage[II]{27}
foregoing therefore proves, taking into account the axiom $(T_0)$, which
holds in $D_+(f)$ and in $\operatorname{Spec}(S_{(f)})$, that $\psi_f$ is
injective and that the inverse map
$\psi_f(D_+(f))\longrightarrow D_+(f)$ is continuous. Finally, to see that
$\psi_f$ is surjective, we remark that if $\mathfrak{q}_0$ is a prime
ideal of $S_{(f)}$ and if, for every $n>0$, we denote by
$\mathfrak{p}_n$ the set of $x\in S_n$ such that
$x^d/f^n\in\mathfrak{q}_0$, the $\mathfrak{p}_n$ satisfy the conditions
of (\hyperref[II.2.1.9]{2.1.9}): indeed, if $x\in S_n$, $y\in S_n$
are such that $x^d/f^n\in\mathfrak{q}_0$ and
$y^d/f^n\in\mathfrak{q}_0$, we have
$(x+y)^{2d}/f^{2n}\in\mathfrak{q}_0$, whence
$(x+y)^d/f^n\in\mathfrak{q}_0$ since $\mathfrak{q}_0$ is prime;
this proves that the $\mathfrak{p}_n$ are subgroups of the $S_n$, and the
verification of the other conditions of
(\hyperref[II.2.1.9]{2.1.9}) is immediate, taking into account that
$\mathfrak{q}_0$ is prime. If $\mathfrak{p}$ is the graded prime ideal
of $S$ thus defined, we indeed have
$\psi_f(\mathfrak{p})=\mathfrak{q}_0$, since for $x\in S_{nd}$, the relations
$x/f^n\in\mathfrak{q}_0$ and $x^d/f^{nd}\in\mathfrak{q}_0$ are
equivalent, $\mathfrak{q}_0$ being prime.
\end{proof}

\begin{corollary}[2.3.7]
\label{II.2.3.7}
In order that $D_+(f)=\varnothing$, it is necessary and sufficient that $f$
be nilpotent.
\end{corollary}

\begin{proof}
Indeed, in order that $\operatorname{Spec}(S_{(f)})=\varnothing$, it is
necessary and sufficient that $S_{(f)}=0$, or again that $1=0$ in $S_f$,
which by definition means that $f$ is nilpotent.
\end{proof}

\begin{corollary}[2.3.8]
\label{II.2.3.8}
Let $E$ be a subset of $S_+$. The following conditions are equivalent:
\begin{enumerate}
  \item[\rm a)] $V_+(E)=X=\operatorname{Proj}(S)$.
  \item[\rm b)] Every element of $E$ is nilpotent.
  \item[\rm c)] The homogeneous components of every element of $E$ are
    nilpotent.
\end{enumerate}
\end{corollary}

\begin{proof}
It is clear that c) implies b) and that b) implies a). If $\mathfrak{J}$ is
the graded ideal of $S$ generated by $E$, condition a) is equivalent to
$V_+(\mathfrak{J})=X$; \emph{a fortiori}, a) implies that every homogeneous
element $f\in\mathfrak{J}$ is such that $V_+(f)=X$, hence $f$ is nilpotent
by virtue of (\hyperref[II.2.3.7]{2.3.7}).
\end{proof}

\begin{corollary}[2.3.9]
\label{II.2.3.9}
If $\mathfrak{J}$ is a graded ideal of $S_+$,
$r_+(\mathfrak{J})$ is the intersection of the graded prime ideals of
$S_+$ which contain $\mathfrak{J}$.
\end{corollary}

\begin{proof}
By considering the graded ring $S/\mathfrak{J}$, we can reduce to the case
where $\mathfrak{J}=0$. We must prove that if $f\in S_+$ is not
nilpotent, there is a graded prime ideal of $S$ not containing $f$; now,
at least one of the homogeneous components of $f$ is not nilpotent, and we
may therefore suppose $f$ homogeneous; the assertion then follows from
(\hyperref[II.2.3.7]{2.3.7}).
\end{proof}

\begin{env}[2.3.10]
\label{II.2.3.10}
For every subset $Y$ of $X=\operatorname{Proj}(S)$, let
$\mathfrak{j}_+(Y)$ be the set of $f\in S_+$ such that
$Y\subset V_+(f)$; it amounts to the same to say that
$\mathfrak{j}_+(Y)=\mathfrak{j}(Y)\cap S_+$; $\mathfrak{j}_+(Y)$ is
therefore an ideal of $S_+$, equal to its radical in $S_+$.
\end{env}

\begin{proposition}[2.3.11]
\label{II.2.3.11}
\begin{enumerate}
  \item[\rm(i)] For every subset $E$ of $S_+$,
    $\mathfrak{j}_+(V_+(E))$ is the radical in $S_+$ of the graded ideal of
    $S_+$ generated by $E$.
  \item[\rm(ii)] For every subset $Y$ of $X$,
    $V_+(\mathfrak{j}_+(Y))=\overline{Y}$, the closure of $Y$ in $X$.
\end{enumerate}
\end{proposition}

\begin{proof}
\begin{enumerate}
  \item[\rm(i)] If $\mathfrak{J}$ is the graded ideal of $S_+$ generated
    by $E$, we have $V_+(E)=V_+(\mathfrak{J})$, and the assertion then
    follows from (\hyperref[II.2.3.9]{2.3.9}).
  \item[\rm(ii)] Since
    $V_+(\mathfrak{J})=\bigcap_{f\in\mathfrak{J}}V_+(f)$, the relation
    $Y\subset V_+(\mathfrak{J})$ implies $Y\subset V_+(f)$ for every
    $f\in\mathfrak{J}$, and consequently
    $\mathfrak{j}_+(Y)\supset\mathfrak{J}$, whence
    $V_+(\mathfrak{j}_+(Y))\subset V_+(\mathfrak{J})$, which proves (ii)
    by definition of the closed sets.
\end{enumerate}
\end{proof}

\begin{corollary}[2.3.12]
\label{II.2.3.12}
The closed subsets $Y$ of $X=\operatorname{Proj}(S)$ and the graded ideals
of $S_+$ equal to their radical in $S_+$ correspond bijectively to each
other under the decreasing maps
$Y\longmapsto\mathfrak{j}_+(Y)$, $\mathfrak{J}\longmapsto
V_+(\mathfrak{J})$; to the union $Y_1\cup Y_2$ of two closed subsets of
$X$ corresponds
\oldpage[II]{28}
$\mathfrak{j}_+(Y_1)\cap\mathfrak{j}_+(Y_2)$, and to the intersection of an
arbitrary family $(Y_\lambda)$ of closed subsets corresponds the radical
in $S_+$ of the sum of the $\mathfrak{j}_+(Y_\lambda)$.
\end{corollary}

\begin{corollary}[2.3.13]
\label{II.2.3.13}
Let $\mathfrak{J}$ be a graded ideal in $S_+$; in order that
$V_+(\mathfrak{J})=\varnothing$, it is necessary and sufficient that every
element of $S_+$ have a power in $\mathfrak{J}$.
\end{corollary}

This last corollary can also be expressed in any of the following
equivalent forms:

\begin{corollary}[2.3.14]
\label{II.2.3.14}
Let $(f_\alpha)$ be a family of homogeneous elements of $S_+$. For the
$D_+(f_\alpha)$ to form a covering of
$X=\operatorname{Proj}(S)$, it is necessary and sufficient that every
element of $S_+$ have a power in the ideal generated by the $f_\alpha$.
\end{corollary}

\begin{corollary}[2.3.15]
\label{II.2.3.15}
Let $(f_\alpha)$ be a family of homogeneous elements of $S_+$, and $f$ an
element of $S_+$. The following relations are equivalent:
\begin{enumerate}
  \item[\rm a)] $D_+(f)\subset\bigcup_\alpha D_+(f_\alpha)$;
  \item[\rm b)] $V_+(f)\supset\bigcap_\alpha V_+(f_\alpha)$;
  \item[\rm c)] some power of $f$ belongs to the ideal generated by the
    $f_\alpha$.
\end{enumerate}
\end{corollary}

\begin{corollary}[2.3.16]
\label{II.2.3.16}
For $X=\operatorname{Proj}(S)$ to be empty, it is necessary and sufficient
that every element of $S_+$ be nilpotent.
\end{corollary}

\begin{corollary}[2.3.17]
\label{II.2.3.17}
In the bijective correspondence described in
(\hyperref[II.2.3.12]{2.3.12}), the irreducible closed subsets of $X$
correspond to the graded prime ideals in $S_+$.
\end{corollary}

\begin{proof}
Indeed, if $Y=Y_1\cup Y_2$, where $Y_1$ and $Y_2$ are closed and distinct
from $Y$, we have
\[
  \mathfrak{j}_+(Y)
  =\mathfrak{j}_+(Y_1)\cap\mathfrak{j}_+(Y_2)
\]
the ideals $\mathfrak{j}_+(Y_1)$ and $\mathfrak{j}_+(Y_2)$ being distinct
from $\mathfrak{j}_+(Y)$, so $\mathfrak{j}_+(Y)$ is not prime.
Conversely, if $\mathfrak{J}$ is a non-prime graded ideal of $S_+$, there
exist two elements $f$, $g$ of $S_+$ such that
$f\notin\mathfrak{J}$, $g\notin\mathfrak{J}$, $fg\in\mathfrak{J}$; then
$V_+(f)\not\supset V_+(\mathfrak{J})$,
$V_+(g)\not\supset V_+(\mathfrak{J})$, but
$V_+(\mathfrak{J})\subset V_+(f)\cup V_+(g)$ by
(\hyperref[II.2.3.2.3]{2.3.2.3}); we conclude that
$V_+(\mathfrak{J})$ is the union of the closed sets
$V_+(f)\cap V_+(\mathfrak{J})$ and $V_+(g)\cap V_+(\mathfrak{J})$, which are
distinct from $V_+(\mathfrak{J})$.
\end{proof}


\subsection{The scheme structure on $\operatorname{Proj}(S)$.}
\label{subsection:II.2.4}

\begin{env}[2.4.1]
\label{II.2.4.1}
Let $f$, $g$ be two homogeneous elements of $S_+$; consider the affine
schemes $Y_f=\operatorname{Spec}(S_{(f)})$,
$Y_g=\operatorname{Spec}(S_{(g)})$ and
$Y_{fg}=\operatorname{Spec}(S_{(fg)})$. By virtue of
(\hyperref[II.2.2.2]{2.2.2}), the morphism
$w_{fg,f}=({}^a\omega_{fg,f},\widetilde{\omega}_{fg,f})$ from $Y_{fg}$ to
$Y_f$ corresponding to the canonical homomorphism
$\omega_{fg,f}:S_{(f)}\to S_{(fg)}$ is an \emph{open immersion}
(\hyperref[I.1.3.6]{I, 1.3.6}). By means of the inverse homeomorphism
of $\psi_f:D_+(f)\to Y_f$ (\hyperref[II.2.3.6]{2.3.6}), we can
transport to $D_+(f)$ the affine scheme structure of $Y_f$; by virtue of
the commutativity of the diagram (\hyperref[II.2.3.5.1]{2.3.5.1}), the
affine scheme $D_+(fg)$ is thereby identified with the scheme induced on the
open set $D_+(fg)$ of the space underlying the affine scheme $D_+(f)$. It is
then clear (taking into account (\hyperref[II.2.3.4]{2.3.4})) that
$X=\operatorname{Proj}(S)$ is endowed with a unique \emph{prescheme}
structure whose restriction to each $D_+(f)$ is the affine scheme just
defined. Moreover:
\end{env}

\begin{proposition}[2.4.2]
\label{II.2.4.2}
The prescheme $\operatorname{Proj}(S)$ is a scheme.
\end{proposition}

\begin{proof}
It suffices (\hyperref[I.5.5.6]{I, 5.5.6}) to show that, for all
homogeneous $f$, $g$ in $S_+$,
$D_+(f)\cap D_+(g)=D_+(fg)$ is affine and its ring is generated by
the canonical images of the rings of $D_+(f)$ and $D_+(g)$; the first
point is evident by definition, and the second follows from
(\hyperref[II.2.2.4]{2.2.4}).
\end{proof}
\oldpage[II]{29}
Whenever we speak of the homogeneous prime spectrum $\operatorname{Proj}(S)$
as a \emph{scheme}, it will always be with the structure just
defined.

\begin{example}[2.4.3]
\label{II.2.4.3}
Take $S=K[T_1,T_2]$, $K$ being a field, $T_1$, $T_2$ two
indeterminates, and $S$ being graded by total degree. It follows from
(\hyperref[II.2.3.14]{2.3.14}) that $\operatorname{Proj}(S)$ is the
union of $D_+(T_1)$ and $D_+(T_2)$; we see immediately that these affine
schemes are canonically isomorphic to $K[T]$, and that
$\operatorname{Proj}(S)$ is obtained by the gluing of these two affine
schemes described in (\hyperref[I.2.3.2]{I, 2.3.2}) (cf.
(\hyperref[II.7.4.14]{7.4.14})).
\end{example}

\begin{proposition}[2.4.4]
\label{II.2.4.4}
Let $S$ be a ring graded in positive degrees, and let $X$ be the scheme
$\operatorname{Proj}(S)$.
\begin{enumerate}
  \item[\rm(i)] If $\mathfrak{N}_+$ is the nilradical of $S_+$
    (\hyperref[II.2.1.10]{2.1.10}), the scheme $X_{\mathrm{red}}$ is
    canonically isomorphic to $\operatorname{Proj}(S/\mathfrak{N}_+)$; in
    particular, if $S$ is essentially reduced,
    $\operatorname{Proj}(S)$ is reduced.
  \item[\rm(ii)] Suppose $S$ essentially reduced. Then, for $X$ to be
    integral, it is necessary and sufficient that $S$ be essentially
    integral.
\end{enumerate}
\end{proposition}

\begin{proof}
\begin{enumerate}
  \item[\rm(i)] Let $\overline{S}$ be the graded ring
    $S/\mathfrak{N}_+$, and denote by $x\mapsto\overline{x}$ the
    canonical homomorphism $S\to\overline{S}$, of degree $0$. For every
    $f\in S_d$ ($d>0$), the canonical homomorphism
    $S_f\to\overline{S}_{\overline{f}}$
    (\hyperref[0.1.5.1]{0, 1.5.1}) is surjective and of degree $0$, and
    therefore gives by restriction a surjective homomorphism
    $S_{(f)}\to\overline{S}_{(\overline{f})}$; if we suppose
    $f\notin\mathfrak{N}_+$, we verify immediately that
    $\overline{S}_{(\overline{f})}$ is reduced and that the kernel of
    the preceding homomorphism is the nilradical of $S_{(f)}$, in other
    words $\overline{S}_{(\overline{f})}=(S_{(f)})_{\mathrm{red}}$. To this
    homomorphism there therefore corresponds a closed immersion
    $D_+(\overline{f})\to D_+(f)$ which identifies $D_+(\overline{f})$ with
    $(D_+(f))_{\mathrm{red}}$ (\hyperref[I.5.1.2]{I, 5.1.2}), and is in
    particular a homeomorphism of the underlying spaces of these two
    affine schemes. Moreover, if $g\notin\mathfrak{N}_+$ is a second
    homogeneous element of $S_+$, the diagram
    \[
      \xymatrix{
        S_{(f)}\ar[r]\ar[d]
        & \overline{S}_{(\overline{f})}\ar[d]\\
        S_{(fg)}\ar[r]
        & \overline{S}_{(\overline{fg})}
      }
    \]
    is commutative; since moreover the $D_+(f)$, for $f$ homogeneous of
    degree $>0$ and $f\notin\mathfrak{N}_+$, form a covering of
    $X=\operatorname{Proj}(S)$ (\hyperref[II.2.3.7]{2.3.7}), we see
    that the morphisms $D_+(\overline{f})\to D_+(f)$ are the restrictions
    of a closed immersion
    $\operatorname{Proj}(\overline{S})\to\operatorname{Proj}(S)$ which is
    a homeomorphism of the underlying spaces; whence the conclusion
    (\hyperref[I.5.1.2]{I, 5.1.2}).
  \item[\rm(ii)] Suppose that $S$ is essentially integral, in other
    words that $(0)$ is a graded prime ideal of $S_+$ distinct from $S_+$;
    then $X$ is reduced by virtue of (i) and irreducible by virtue of
    (\hyperref[II.2.3.17]{2.3.17}). Conversely, suppose that $S$ is
    essentially reduced and $X$ integral; then, for $f\neq0$ homogeneous
    in $S_+$, we have $D_+(f)\neq\varnothing$
    (\hyperref[II.2.3.7]{2.3.7}); the hypothesis that $X$ is irreducible
    implies that $D_+(f)\cap D_+(g)\neq\varnothing$ for $f$, $g$
    homogeneous and $\neq0$ in $S_+$; hence $fg\neq0$ by virtue of
    (\hyperref[II.2.3.3.2]{2.3.3.2}), and we conclude that $S_+$ has no
    zero divisors, whence the first assertion.
\end{enumerate}
\end{proof}

\begin{env}[2.4.5]
\label{II.2.4.5}
Given a commutative ring $A$, recall that a graded ring $S$ is said to be
a \emph{graded $A$-algebra} if it is endowed with an $A$-algebra structure
such that each of the subgroups $S_n$ is an $A$-module;
moreover, for this it suffices that $S_0$ be
\oldpage[II]{30}
an $A$-algebra; in other words, we define a graded $A$-algebra structure on
$S$ by defining an $A$-algebra structure on $S_0$
and setting, for $\alpha\in A$, $x\in S_n$,
$\alpha\cdot x=(\alpha\cdot1)x$.
\end{env}

\begin{proposition}[2.4.6]
\label{II.2.4.6}
Suppose that $S$ is a graded $A$-algebra. Then on
$X=\operatorname{Proj}(S)$, the structural sheaf $\mathcal{O}_X$ is an
$A$-algebra ($A$ being regarded as a simple sheaf on $X$); in
other words, $X$ is a scheme over $\operatorname{Spec}(A)$.
\end{proposition}

\begin{proof}
It suffices to note that for every homogeneous $f$ in $S_+$, $S_{(f)}$ is
an algebra over $A$, and that the diagram
\[
  \xymatrix{
    S_{(f)}\ar[rr] && S_{(fg)}\\
    & A\ar[ul]\ar[ur] &
  }
\]
is commutative, for homogeneous $f$, $g$ in $S_+$.
\end{proof}

\begin{proposition}[2.4.7]
\label{II.2.4.7}
Let $S$ be a ring graded in positive degrees.
\begin{enumerate}
  \item[\rm(i)] For every integer $d>0$, there exists a canonical isomorphism
    of the scheme $\operatorname{Proj}(S)$ onto the scheme
    $\operatorname{Proj}(S^{(d)})$.
  \item[\rm(ii)] Let $S'$ be the graded ring such that $S'_0=\mathbf{Z}$,
    $S'_n=S_n$ (regarded as a $\mathbf{Z}$-module) for $n>0$. There exists
    a canonical isomorphism of the scheme $\operatorname{Proj}(S)$ onto the
    scheme $\operatorname{Proj}(S')$.
\end{enumerate}
\end{proposition}

\begin{proof}
\begin{enumerate}
  \item[\rm(i)] Let us first show that the map
    $\mathfrak{p}\mapsto\mathfrak{p}\cap S^{(d)}$ is a bijection of
    the set $\operatorname{Proj}(S)$ onto
    $\operatorname{Proj}(S^{(d)})$. Indeed, suppose given a graded
    prime ideal $\mathfrak{p}'\in\operatorname{Proj}(S^{(d)})$, and set
    $\mathfrak{p}_{nd}=\mathfrak{p}'\cap S_{nd}$ ($n\geq0$). For every
    $n>0$ which is not a multiple of $d$, define $\mathfrak{p}_n$ as
    the set of $x\in S_n$ such that
    $x^d\in\mathfrak{p}_{nd}$; if $x\in\mathfrak{p}_n$,
    $y\in\mathfrak{p}_n$, we have
    $(x+y)^{2d}\in\mathfrak{p}_{2nd}$, hence
    $(x+y)^d\in\mathfrak{p}_{nd}$ since $\mathfrak{p}'$ is prime; it
    is immediate that the $\mathfrak{p}_n$ thus defined for every $n\geq0$
    satisfy the conditions of
    (\hyperref[II.2.1.9]{2.1.9}), so there exists a unique prime ideal
    $\mathfrak{p}\in\operatorname{Proj}(S)$ such that
    $\mathfrak{p}\cap S^{(d)}=\mathfrak{p}'$. Since for every homogeneous
    $f$ in $S_+$ we have
    $V_+(f)=V_+(f^d)$ (\hyperref[II.2.3.2.3]{2.3.2.3}), we see that the
    preceding bijection is a homeomorphism of topological spaces.
    Finally, with the same notation, $S_{(f)}$ and
    $S^{(d)}_{(f^d)}$ are canonically identified
    (\hyperref[II.2.2.2]{2.2.2}), so $\operatorname{Proj}(S)$ and
    $\operatorname{Proj}(S^{(d)})$ are canonically identified as
    schemes.
  \item[\rm(ii)] If to every
    $\mathfrak{p}\in\operatorname{Proj}(S)$ we assign the unique
    prime ideal $\mathfrak{p}'\in\operatorname{Proj}(S')$ such that
    $\mathfrak{p}'\cap S_n=\mathfrak{p}\cap S_n$ for every $n>0$
    (\hyperref[II.2.1.9]{2.1.9}), it is clear that we define a
    canonical homeomorphism
    $\operatorname{Proj}(S)\xrightarrow{\sim}\operatorname{Proj}(S')$ of the
    underlying spaces, since $V_+(f)$ is the same set for $S$ and
    $S'$ when $f$ is a homogeneous element of $S_+$. Since moreover
    $S_{(f)}=S'_{(f)}$, $\operatorname{Proj}(S)$ and
    $\operatorname{Proj}(S')$ are identified as schemes.
\end{enumerate}
\end{proof}

\begin{corollary}[2.4.8]
\label{II.2.4.8}
If $S$ is a graded $A$-algebra and $S'_A$ the graded $A$-algebra such that
$(S'_A)_0=A$, $(S'_A)_n=S_n$ for $n>0$, there exists a canonical isomorphism
of $\operatorname{Proj}(S)$ onto $\operatorname{Proj}(S'_A)$.
\end{corollary}

\begin{proof}
Indeed, these two schemes are canonically isomorphic to
$\operatorname{Proj}(S')$, with the notation of
(\hyperref[II.2.4.7]{2.4.7, (ii)}).
\end{proof}

\subsection{Sheaf associated to a graded module.}
\label{subsection:II.2.5}

\begin{env}[2.5.1]
\label{II.2.5.1}
Let $M$ be a graded $S$-module. For every homogeneous $f$ in $S_+$,
$M_{(f)}$ is an $S_{(f)}$-module, and there therefore corresponds to it an
associated quasi-coherent sheaf $\widetilde{M_{(f)}}$ on the affine scheme
$\operatorname{Spec}(S_{(f)})$, identified with $D_+(f)$
(\hyperref[I.1.3.4]{I, 1.3.4}).
\end{env}
\oldpage[II]{31}

\begin{proposition}[2.5.2]
\label{II.2.5.2}
There exists on $X=\operatorname{Proj}(S)$ one and only one quasi-coherent
$\mathcal{O}_X$-module $\widetilde{M}$ such that for every homogeneous $f$ in
$S_+$ we have
$\Gamma(D_+(f),\widetilde{M})=M_{(f)}$, the restriction homomorphism
$\Gamma(D_+(f),\widetilde{M})\to
\Gamma(D_+(fg),\widetilde{M})$, for homogeneous $f$, $g$ in $S_+$, being
the canonical homomorphism $M_{(f)}\to M_{(fg)}$
(\hyperref[II.2.2.3]{2.2.3}).
\end{proposition}

\begin{proof}
Suppose $f\in S_d$, $g\in S_e$. Since $D_+(fg)$ is identified with the
prime spectrum of $(S_{(f)})_{g^d/f^e}$ by virtue of
(\hyperref[II.2.2.2]{2.2.2}), the restriction to $D_+(fg)$ of the sheaf
$\widetilde{M_{(f)}}$ on $D_+(f)$ is canonically identified with the sheaf
associated to the module $(M_{(f)})_{g^d/f^e}$
(\hyperref[I.1.3.6]{I, 1.3.6}), hence also with
$\widetilde{M_{(fg)}}$ (\hyperref[II.2.2.2]{2.2.2}); we conclude
that there exists a canonical isomorphism
\[
  \theta_{g,f}:\widetilde{M_{(f)}}|D_+(fg)
  \xrightarrow{\sim}\widetilde{M_{(g)}}|D_+(fg)
\]
such that, if $h$ is a third homogeneous element of $S_+$, we have
$\theta_{f,h}=\theta_{f,g}\circ\theta_{g,h}$ in $D_+(fgh)$. Consequently
(\hyperref[0.3.3.1]{0, 3.3.1}) there exist on $X$ a quasi-coherent
$\mathcal{O}_X$-module $\mathcal{F}$ and, for each homogeneous $f$
in $S_+$, an isomorphism $\eta_f$ of
$\mathcal{F}|D_+(f)$ onto $\widetilde{M_{(f)}}$, such that
$\theta_{g,f}=\eta_g\circ\eta_f^{-1}$. If we then consider the sheaf
$\mathcal{G}$ associated to the presheaf (on the basis of the topology of
$X$ formed by the $D_+(f)$) defined by $D_+(f)\to M_{(f)}$, with the
canonical homomorphisms $M_{(f)}\to M_{(fg)}$ as restriction
homomorphisms, the foregoing proves that $\mathcal{F}$ and $\mathcal{G}$
are isomorphic (taking into account (\hyperref[I.1.3.7]{I, 1.3.7})); the
sheaf $\mathcal{G}$ will be denoted by $\widetilde{M}$, and it indeed
satisfies the conditions of the statement. In particular, we have
$\widetilde{S}=\mathcal{O}_X$.
\end{proof}

\begin{definition}[2.5.3]
\label{II.2.5.3}
We say that the quasi-coherent $\mathcal{O}_X$-module $\widetilde{M}$
defined in (\hyperref[II.2.5.2]{2.5.2}) is \emph{associated} to the
graded $S$-module $M$.
\end{definition}

Recall that the graded $S$-modules form a category when the
homomorphisms of graded modules are restricted to the homomorphisms
\emph{of degree $0$}. With this convention:

\begin{proposition}[2.5.4]
\label{II.2.5.4}
The functor $M\mapsto\widetilde{M}$ is an exact additive covariant functor
from the category of graded $S$-modules to that of quasi-coherent
$\mathcal{O}_X$-modules, which commutes with inductive limits and direct
sums.
\end{proposition}

\begin{proof}
Indeed, since these properties are local, it suffices to verify them on
the sheaves
$\widetilde{M}|D_+(f)=\widetilde{M_{(f)}}$; now the functor
$M\mapsto M_f$, the functor $N\mapsto N_0$ (on the category of graded
$S_f$-modules) and the functor $P\mapsto\widetilde{P}$ (on the
category of $S_{(f)}$-modules) all three have the properties of exactness
and of commutation with inductive limits and direct sums
(\hyperref[I.1.3.5]{I, 1.3.5} and
\hyperref[I.1.3.9]{I, 1.3.9}); whence the proposition.
\end{proof}

We denote by $\widetilde{u}$ the homomorphism
$\widetilde{M}\to\widetilde{N}$ corresponding to a homomorphism of degree
$0$, $u:M\to N$. We deduce immediately from
(\hyperref[II.2.5.4]{2.5.4}) that the results of
(\hyperref[I.1.3.9]{I, 1.3.9} and
\hyperref[I.1.3.10]{I, 1.3.10}) remain valid for graded
$S$-modules and homomorphisms of degree $0$ (with the meaning given
here to $\widetilde{M}$), the proofs being purely formal.

\begin{proposition}[2.5.5]
\label{II.2.5.5}
For every $\mathfrak{p}\in X=\operatorname{Proj}(S)$, we have
$\widetilde{M}_{\mathfrak{p}}=M_{(\mathfrak{p})}$.
\end{proposition}

\begin{proof}
Indeed, by definition
$\widetilde{M}_{\mathfrak{p}}=
\varinjlim\Gamma(D_+(f),\widetilde{M})$, where $f$ runs through the set of
homogeneous elements of $S_+$ such that $f\notin\mathfrak{p}$; since
$\Gamma(D_+(f),\widetilde{M})=M_{(f)}$, the proposition follows from the
definition of $M_{(\mathfrak{p})}$
(\hyperref[II.2.2.7]{2.2.7}).
\end{proof}
\oldpage[II]{32}

In particular, \emph{the local ring $(\mathcal{O}_X)_\mathfrak{p}$} is none
other than the ring $S_{(\mathfrak{p})}$, the set of fractions $x/f$, where
$f$ is homogeneous in $S_+$ and does not belong to $\mathfrak{p}$, and $x$
is homogeneous of the same degree as $f$.

More particularly, if $S$ is \emph{essentially integral}, so
that $\operatorname{Proj}(S)=X$ is \emph{integral}
(\hyperref[II.2.4.4]{2.4.4}), and if $\xi=(0)$ is the generic point of
$X$, the \emph{field of rational functions}
$R(X)=\mathcal{O}_\xi$ is the field consisting of the elements $fg^{-1}$,
where $f$ and $g$ are homogeneous of the same degree in $S_+$ and $g\neq0$.

\begin{proposition}[2.5.6]
\label{II.2.5.6}
If, for every $z\in M$ and every homogeneous $f$ in $S_+$, there exists a
power of $f$ annihilating $z$, then $\widetilde{M}=0$. This sufficient
condition is also necessary when the $S_0$-algebra $S$ is generated
by the set $S_1$ of homogeneous elements of degree $1$.
\end{proposition}

\begin{proof}
Indeed, the condition $\widetilde{M}=0$ is equivalent to $M_{(f)}=0$ for
every homogeneous $f$ in $S_+$. On the other hand, if $f\in S_d$, to say that
$M_{(f)}=0$ means that for every $z\in M$ homogeneous of degree a multiple
of $d$, there is a power $f^n$ such that $f^nz=0$. If $M_{(f)}=0$ for every
$f\in S_1$, the condition of the statement is therefore satisfied for every
$f\in S_1$; it is satisfied \emph{a fortiori} for every homogeneous $f$ in
$S_+$ when $S_1$ generates $S$, since every homogeneous element of $S_+$
is then a linear combination of products of elements of $S_1$.
\end{proof}

\begin{proposition}[2.5.7]
\label{II.2.5.7}
Let $d$ be an integer $>0$, and $f\in S_d$. Then, for every
$n\in\mathbf{Z}$, the $(\mathcal{O}_X|D_+(f))$-module
$\widetilde{S(nd)}|D_+(f)$ is canonically isomorphic to
$\mathcal{O}_X|D_+(f)$.
\end{proposition}

\begin{proof}
Indeed, multiplication by the invertible element $(f/1)^n$ of $S_f$
is a bijection of $S_{(f)}=(S_f)_0$ onto
$(S_f)_{nd}=(S_f(nd))_0=(S(nd)_f)_0=S(nd)_{(f)}$; in other words, the
$S_{(f)}$-modules $S_{(f)}$ and $S(nd)_{(f)}$ are canonically
isomorphic, whence the proposition.
\end{proof}

\begin{corollary}[2.5.8]
\label{II.2.5.8}
On the open set
\[
  U=\bigcup_{f\in S_d}D_+(f),
\]
the restriction of the $\mathcal{O}_X$-module $\widetilde{S(nd)}$ is an
invertible $(\mathcal{O}_X|U)$-module
(\hyperref[0.5.4.1]{0, 5.4.1}).
\end{corollary}

\begin{corollary}[2.5.9]
\label{II.2.5.9}
If the ideal $S_+$ of $S$ is generated by the set $S_1$ of homogeneous
elements of degree $1$, the $\mathcal{O}_X$-modules
$\widetilde{S(n)}$ are invertible for every $n\in\mathbf{Z}$.
\end{corollary}

\begin{proof}
It suffices to note that
$X=\bigcup_{f\in S_1}D_+(f)$ by virtue of the hypothesis
(\hyperref[II.2.3.14]{2.3.14}) and to apply
(\hyperref[II.2.5.8]{2.5.8}) with $U=X$.
\end{proof}

\begin{env}[2.5.10]
\label{II.2.5.10}
Throughout the rest of this section, we set
\[
\label{II.2.5.10.1}
  \mathcal{O}_X(n)=\widetilde{S(n)}
\tag{2.5.10.1}
\]
for every $n\in\mathbf{Z}$, and, for every open subset $U$ of $X$ and every
$(\mathcal{O}_X|U)$-module $\mathcal{F}$,
\[
\label{II.2.5.10.2}
  \mathcal{F}(n)=
  \mathcal{F}\otimes_{\mathcal{O}_X|U}(\mathcal{O}_X(n)|U)
\tag{2.5.10.2}
\]
for every $n\in\mathbf{Z}$. If the ideal $S_+$ is generated by $S_1$, the
functor $\mathcal{F}(n)$ is \emph{exact} in
$\mathcal{F}$ for every $n\in\mathbf{Z}$, since $\mathcal{O}_X(n)$ is then
an invertible $\mathcal{O}_X$-module.
\end{env}

\begin{env}[2.5.11]
\label{II.2.5.11}
Let $M$ and $N$ be two graded $S$-modules. For every $f\in S_d$ ($d>0$)
we define a canonical functorial homomorphism of $S_{(f)}$-modules
\[
\label{II.2.5.11.1}
  \lambda_f:M_{(f)}\otimes_{S_{(f)}}N_{(f)}
  \longrightarrow(M\otimes_S N)_{(f)}
\tag{2.5.11.1}
\]
\oldpage[II]{33}

by composing the homomorphism
$M_{(f)}\otimes_{S_{(f)}}N_{(f)}\to M_f\otimes_{S_f}N_f$
(coming from the injections $M_{(f)}\to M_f$, $N_{(f)}\to N_f$ and
$S_{(f)}\to S_f$) with the canonical isomorphism
$M_f\otimes_{S_f}N_f\xrightarrow{\sim}(M\otimes_S N)_f$
(\hyperref[0.1.3.4]{0, 1.3.4}), and noting that, by definition of the
tensor product of two graded modules, the latter isomorphism preserves
degrees; for $x\in M_{md}$, $y\in N_{nd}$ ($m\geq0$, $n\geq0$), we
therefore have
\[
  \lambda_f((x/f^m)\otimes(y/f^n))=(x\otimes y)/f^{m+n}.
\]

It follows immediately from this definition that if $g\in S_e$ ($e>0$), the
diagram
\[
  \xymatrix{
    M_{(f)}\otimes_{S_{(f)}}N_{(f)}\ar[r]^{\lambda_f}\ar[d]
    & (M\otimes_S N)_{(f)}\ar[d]\\
    M_{(fg)}\otimes_{S_{(fg)}}N_{(fg)}\ar[r]_{\lambda_{fg}}
    & (M\otimes_S N)_{(fg)}
  }
\]
(where the right vertical arrow is the canonical homomorphism and the left
one comes from the canonical homomorphisms) is commutative. We therefore
deduce from the $\lambda$ a canonical functorial homomorphism of
$\mathcal{O}_X$-modules
\[
\label{II.2.5.11.2}
  \lambda:\widetilde{M}\otimes_{\mathcal{O}_X}\widetilde{N}
  \longrightarrow\widetilde{M\otimes_S N}.
\tag{2.5.11.2}
\]

Consider in particular two graded ideals $\mathfrak{J}$,
$\mathfrak{K}$ of $S$; since $\widetilde{\mathfrak{J}}$ and
$\widetilde{\mathfrak{K}}$ are sheaves of ideals of
$\mathcal{O}_X$, we have a canonical homomorphism
$\widetilde{\mathfrak{J}}\otimes_{\mathcal{O}_X}
\widetilde{\mathfrak{K}}\to\mathcal{O}_X$; the diagram
\[
\label{II.2.5.11.3}
  \xymatrix{
    \widetilde{\mathfrak{J}}\otimes_{\mathcal{O}_X}
    \widetilde{\mathfrak{K}}\ar[rr]^\lambda\ar[dr]
    && \widetilde{\mathfrak{J}\otimes_S\mathfrak{K}}\ar[dl]\\
    & \mathcal{O}_X &
  }
\tag{2.5.11.3}
\]
is then commutative. Indeed, we are reduced to verifying this in each
open subset $D_+(f)$ ($f$ homogeneous in $S_+$), and it follows immediately
from the definition (\hyperref[II.2.5.11.1]{2.5.11.1}) of $\lambda_f$ and
from (\hyperref[I.1.3.13]{I, 1.3.13}).

Finally, note that if $M$, $N$, $P$ are three graded $S$-modules, the
diagram
\[
\label{II.2.5.11.4}
  \xymatrix{
    \widetilde{M}\otimes_{\mathcal{O}_X}\widetilde{N}
    \otimes_{\mathcal{O}_X}\widetilde{P}
    \ar[r]^{\lambda\otimes1}\ar[d]_{1\otimes\lambda}
    & \widetilde{M\otimes_S N}\otimes_{\mathcal{O}_X}\widetilde{P}
    \ar[d]^\lambda\\
    \widetilde{M}\otimes_{\mathcal{O}_X}\widetilde{N\otimes_S P}
    \ar[r]_\lambda
    & \widetilde{M\otimes_S N\otimes_S P}
  }
\tag{2.5.11.4}
\]
\oldpage[II]{34}

is commutative. Again it suffices to verify this in each $D_+(f)$, and
it follows immediately from the definitions and from
(\hyperref[I.1.3.13]{I, 1.3.13}).
\end{env}

\begin{env}[2.5.12]
\label{II.2.5.12}
Under the hypotheses of (\hyperref[II.2.5.11]{2.5.11}), we define a
canonical functorial homomorphism of $S_{(f)}$-modules
\[
\label{II.2.5.12.1}
  \mu_f:\bigl(\operatorname{Hom}_S(M,N)\bigr)_{(f)}
  \longrightarrow
  \operatorname{Hom}_{S_{(f)}}(M_{(f)},N_{(f)})
\tag{2.5.12.1}
\]
by assigning to $u/f^n$, where $u$ is a homomorphism of degree
$nd$, the homomorphism $M_{(f)}\to N_{(f)}$ which maps $x/f^m$
($x\in M_{md}$) to $u(x)/f^{m+n}$. For $g\in S_e$ ($e>0$), we again have
a commutative diagram
\[
  \xymatrix{
    \bigl(\operatorname{Hom}_S(M,N)\bigr)_{(f)}
    \ar[r]^{\mu_f}\ar[d]
    & \operatorname{Hom}_{S_{(f)}}(M_{(f)},N_{(f)})\ar[d]\\
    \bigl(\operatorname{Hom}_S(M,N)\bigr)_{(fg)}
    \ar[r]_{\mu_{fg}}
    & \operatorname{Hom}_{S_{(fg)}}(M_{(fg)},N_{(fg)})
  }
\]
(the left arrow being the canonical homomorphism, the right one
coming from the canonical homomorphisms). We again conclude (taking into
account (\hyperref[I.1.3.8]{I, 1.3.8})) that the $\mu_f$ define a
canonical functorial homomorphism of $\mathcal{O}_X$-modules
\[
\label{II.2.5.12.2}
  \mu:\widetilde{\operatorname{Hom}_S(M,N)}
  \longrightarrow
  \mathcal{H}om_{\mathcal{O}_X}(\widetilde{M},\widetilde{N}).
\tag{2.5.12.2}
\]
\end{env}

\begin{proposition}[2.5.13]
\label{II.2.5.13}
Suppose the ideal $S_+$ is generated by $S_1$. Then the homomorphism
$\lambda$ (\hyperref[II.2.5.11.2]{2.5.11.2}) is an isomorphism; so is
the homomorphism $\mu$
(\hyperref[II.2.5.12.2]{2.5.12.2}) when the graded $S$-module $M$
admits a finite presentation (\hyperref[II.2.1.1]{2.1.1}).
\end{proposition}

\begin{proof}
Since $X$ is the union of the $D_+(f)$ for $f\in S_1$
(\hyperref[II.2.3.14]{2.3.14}), we are reduced to proving that
$\lambda_f$ and $\mu_f$ are isomorphisms, under the hypotheses considered,
when $f$ is homogeneous and \emph{of degree $1$}. Now, we then define a
$\mathbf{Z}$-bilinear map
\[
  M_m\times N_n\longrightarrow
  M_{(f)}\otimes_{S_{(f)}}N_{(f)}
\]
by assigning to $(x,y)$ the element
$(x/f^m)\otimes(y/f^n)$ (when $m<0$, we write $x/f^m$ for
$f^{-m}x/1$); these maps define a
$\mathbf{Z}$-linear map
\[
  M\otimes_{\mathbf{Z}}N\longrightarrow
  M_{(f)}\otimes_{S_{(f)}}N_{(f)},
\]
and if $s\in S_q$, this map sends $(sx)\otimes y$ to
\[
  (s/f^q)\bigl((x/f^m)\otimes(y/f^n)\bigr)
\]
(for $x\in M_m$, $y\in N_n$). We therefore deduce a di-homomorphism of
modules
\[
  \gamma_f:M\otimes_S N\longrightarrow
  M_{(f)}\otimes_{S_{(f)}}N_{(f)},
\]
relative to the canonical homomorphism $S\to S_{(f)}$ (sending
$s\in S_q$ to $s/f^q$). Suppose moreover that for an element
$\sum_i(x_i\otimes y_i)$ of $M\otimes_S N$ ($x_i$, $y_i$ homogeneous of
respective degrees $m_i$, $n_i$) we have
\[
  f^r\sum_i(x_i\otimes y_i)=0,
\]
in other words $\sum_i(f^rx_i\otimes y_i)=0$. We deduce by virtue of
(\hyperref[0.1.3.4]{0, 1.3.4}) that
\[
  \sum_i(f^rx_i/f^{m_i+r})\otimes(y_i/f^{n_i})=0,
\]
that is, $\gamma_f(\sum_i(x_i\otimes y_i))=0$. Consequently $\gamma_f$
factors as
\[
  M\otimes_S N\longrightarrow(M\otimes_S N)_f
  \xrightarrow{\gamma'_f}
  M_{(f)}\otimes_{S_{(f)}}N_{(f)}~;
\]
if $\lambda'_f$ is the restriction of $\gamma'_f$
\oldpage[II]{35}
to $(M\otimes_S N)_{(f)}$, we verify immediately that $\lambda_f$ and
$\lambda'_f$ are mutually inverse $S_{(f)}$-homomorphisms, whence the
first part of the proposition.

To prove the second part, suppose that $M$ is the cokernel of a
homomorphism $P\to Q$ of graded $S$-modules, $P$ and $Q$ being direct
sums of finitely many modules of the form $S(n)$; using
the left exactness of $\operatorname{Hom}_S(L,N)$ in $L$ and the exactness
of $M_{(f)}$ in $M$, we are immediately reduced to proving that $\mu_f$ is
an isomorphism when $M=S(n)$. Now, for every homogeneous $z$ in $N$, let
$u_z$ be the homomorphism from $S(n)$ to $N$ such that $u_z(1)=z$; we see
immediately that $\eta:z\mapsto u_z$ is an isomorphism of degree $0$ of
$N(-n)$ onto $\operatorname{Hom}_S(S(n),N)$. To it there corresponds an
isomorphism
\[
  \eta_f:\bigl(N(-n)\bigr)_{(f)}
  \longrightarrow
  \bigl(\operatorname{Hom}_S(S(n),N)\bigr)_{(f)}.
\]
On the other hand, let $\eta'_f$ be the isomorphism
\[
  N_{(f)}\longrightarrow
  \operatorname{Hom}_{S_{(f)}}\bigl(S(n)_{(f)},N_{(f)}\bigr)
\]
which assigns to every $z'\in N_{(f)}$ the homomorphism $v_{z'}$ such
that $v_{z'}(s/f^k)=sz'/f^{n+k}$ (for
$s\in S_{n+k}=(S(n))_k$). We easily check that the composite map
\[
  \bigl(N(-n)\bigr)_{(f)}
  \xrightarrow{\eta_f}
  \bigl(\operatorname{Hom}_S(S(n),N)\bigr)_{(f)}
  \xrightarrow{\mu_f}
  \operatorname{Hom}_{S_{(f)}}\bigl(S(n)_{(f)},N_{(f)}\bigr)
  \xrightarrow{{\eta'_f}^{-1}}N_{(f)}
\]
is the isomorphism $z/f^h\mapsto z/f^{h-n}$ of
$\bigl(N(-n)\bigr)_{(f)}$ onto $N_{(f)}$, so $\mu_f$ is an isomorphism.
\end{proof}

If the ideal $S_+$ is generated by $S_1$, we deduce from
(\hyperref[II.2.5.13]{2.5.13}) that for every graded ideal $\mathfrak{J}$
of $S$ and every graded $S$-module $M$, we have
\[
\label{II.2.5.13.1}
  \widetilde{\mathfrak{J}}\cdot\widetilde{M}
  =\widetilde{\mathfrak{J}\cdot M}
\tag{2.5.13.1}
\]
up to canonical isomorphism; indeed, this follows from the commutativity
of the diagram
\[
  \xymatrix{
    \widetilde{\mathfrak{J}}\otimes_{\mathcal{O}_X}\widetilde{M}
    \ar[rr]^\lambda\ar[dr]
    && \widetilde{\mathfrak{J}\otimes_S M}\ar[dl]\\
    & \widetilde{M} &
  }
\]
which is verified like that of
(\hyperref[II.2.5.11.3]{2.5.11.3}).

\begin{corollary}[2.5.14]
\label{II.2.5.14}
Suppose $S$ is generated by $S_1$. Then, for all $m$, $n$ in
$\mathbf{Z}$, we have:
\[
\label{II.2.5.14.1}
  \mathcal{O}_X(m)\otimes_{\mathcal{O}_X}\mathcal{O}_X(n)
  =\mathcal{O}_X(m+n)
\tag{2.5.14.1}
\]
\[
\label{II.2.5.14.2}
  \mathcal{O}_X(n)=\bigl(\mathcal{O}_X(1)\bigr)^{\otimes n}
\tag{2.5.14.2}
\]
up to canonical isomorphisms.
\end{corollary}

\begin{proof}
The first formula follows from (\hyperref[II.2.5.13]{2.5.13}) and from
the existence of the canonical isomorphism of degree $0$,
$S(m)\otimes_S S(n)\xrightarrow{\sim}S(m+n)$, which maps the element
$1\otimes 1$ (where the first factor $1$ is in $(S(m))_{-m}$ and the second
in $(S(n))_{-n}$) to the element
$1\in(S(m+n))_{-m-n}$. It then suffices to prove the second formula
for $n=-1$, and by virtue of (\hyperref[II.2.5.13]{2.5.13}), this amounts
to seeing that $\operatorname{Hom}_S(S(1),S)$ is canonically isomorphic to
$S(-1)$, an immediate verification by going back to the definitions
(\hyperref[II.2.1.2]{2.1.2}) and remembering that $S(1)$ is a
monogenic $S$-module.
\end{proof}

\oldpage[II]{36}
\begin{corollary}[2.5.15]
\label{II.2.5.15}
Suppose $S$ is generated by $S_1$. For every graded $S$-module $M$ and
every $n\in\mathbf{Z}$, we have
\[
\label{II.2.5.15.1}
  \widetilde{M(n)}=\widetilde{M}(n)
\tag{2.5.15.1}
\]
up to canonical isomorphism.
\end{corollary}

\begin{proof}
This follows from the definitions (\hyperref[II.2.5.10.2]{2.5.10.2}) and
(\hyperref[II.2.5.10.1]{2.5.10.1}), from Proposition
(\hyperref[II.2.5.13]{2.5.13}) and from the existence of a canonical
isomorphism $M(n)\xrightarrow{\sim}M\otimes_S S(n)$ of degree $0$ which maps
every $z\in(M(n))_h=M_{n+h}$ to
$z\otimes1\in M_{n+h}\otimes(S(n))_{-n}\subset(M\otimes_S S(n))_h$.
\end{proof}

\begin{env}[2.5.16]
\label{II.2.5.16}
Denote by $S'$ the graded ring such that $S'_0=\mathbf{Z}$,
$S'_n=S_n$ for $n>0$. Then, if $f\in S_d$ ($d>0$), we have
$(S(n))_{(f)}=(S'(n))_{(f)}$ for every $n\in\mathbf{Z}$, since an element of
$(S'(n))_{(f)}$ is of the form $x/f^k$ with $x\in S'_{n+kd}$ ($k>0$), and
we can always take $k$ such that $n+kd\neq0$. Since
$X=\operatorname{Proj}(S)$ and $X'=\operatorname{Proj}(S')$ are
canonically identified (\hyperref[II.2.4.7]{2.4.7, (ii)}), we see that for
every $n\in\mathbf{Z}$, $\mathcal{O}_X(n)$ and $\mathcal{O}_{X'}(n)$ are
images of each other under the preceding identification.

Note on the other hand that for every $d>0$ and every $n\in\mathbf{Z}$, we
have
\[
  (S^{(d)}(n))_h=S_{(n+h)d}=(S(nd))_{hd}
\]
for $f\in S_d$, we therefore have $(S^{(d)}(n))_{(f)}=(S(nd))_{(f)}$. We
know that the schemes $X=\operatorname{Proj}(S)$ and
$X^{(d)}=\operatorname{Proj}(S^{(d)})$ are canonically identified
(\hyperref[II.2.4.7]{2.4.7, (i)}); the foregoing shows that if the
$S_0$-algebra $S^{(d)}$ is generated by $S_d$, then $\mathcal{O}_X(nd)$ and
$\mathcal{O}_{X^{(d)}}(n)$ are images of each other under this
identification, for every $n\in\mathbf{Z}$.
\end{env}

\begin{proposition}[2.5.17]
\label{II.2.5.17}
Let $d$ be an integer $>0$, and let
$U=\bigcup_{f\in S_d}D_+(f)$. The restriction to $U$ of the canonical
homomorphism
$\mathcal{O}_X(nd)\otimes_{\mathcal{O}_X}\mathcal{O}_X(-nd)
\to\mathcal{O}_X$ is an isomorphism for every integer $n$.
\end{proposition}

\begin{proof}
By virtue of (\hyperref[II.2.5.16]{2.5.16}), we may restrict to the case
where $d=1$, and the conclusion then follows from the proof of
(\hyperref[II.2.5.13]{2.5.13}).
\end{proof}

\subsection{Graded $S$-module associated to a sheaf on
$\operatorname{Proj}(S)$.}
\label{subsection:II.2.6}

\emph{Throughout this number, we suppose that the ideal $S_+$ of $S$ is
generated by the set $S_1$ of homogeneous elements of degree $1$.}

\begin{env}[2.6.1]
\label{II.2.6.1}
The $\mathcal{O}_X$-module $\mathcal{O}_X(1)$ is then
\emph{invertible} (\hyperref[II.2.5.9]{2.5.9}); we then set, for every
$\mathcal{O}_X$-module $\mathcal{F}$
(\hyperref[0.5.4.6]{0, 5.4.6}),
\[
\label{II.2.6.1.1}
  \Gamma_*(\mathcal{F})
  =\Gamma_*(\mathcal{O}_X(1),\mathcal{F})
  =\bigoplus_{n\in\mathbf{Z}}\Gamma(X,\mathcal{F}(n))
\tag{2.6.1.1}
\]
taking into account (\hyperref[II.2.5.14.2]{2.5.14.2}). Recall
(\hyperref[0.5.4.6]{0, 5.4.6}) that $\Gamma_*(\mathcal{O}_X)$ is endowed
with a structure of \emph{graded ring}, and $\Gamma_*(\mathcal{F})$ with a
structure of
\emph{graded module} over $\Gamma_*(\mathcal{O}_X)$.

Since $\mathcal{O}_X(n)$ is locally free, $\Gamma_*(\mathcal{F})$
is an additive covariant functor in $\mathcal{F}$ which is
\emph{left exact}; in particular, if $\mathcal{J}$
is a sheaf of ideals of $\mathcal{O}_X$, $\Gamma_*(\mathcal{J})$
is canonically identified with a
\emph{graded ideal} of $\Gamma_*(\mathcal{O}_X)$.
\end{env}

\begin{env}[2.6.2]
\label{II.2.6.2}
Let $M$ be a graded $S$-module; for every $f\in S_d$ ($d>0$),
$x\mapsto x/1$ is a homomorphism of abelian groups
$M_0\to M_{(f)}$, and since $M_{(f)}$ is canonically

\oldpage[II]{37}

identified with $\Gamma(D_+(f),\widetilde{M})$, we thus obtain a
homomorphism of abelian groups
\[
  \alpha_0^f:M_0\longrightarrow\Gamma(D_+(f),\widetilde{M}).
\]
It is clear that, for every $g\in S_e$ ($e>0$), the diagram
\[
  \xymatrix{
    & \Gamma(D_+(f),\widetilde{M})\ar[dd]\\
    M_0\ar[ur]^{\alpha_0^f}\ar[dr]_{\alpha_0^{fg}}\\
    & \Gamma(D_+(fg),\widetilde{M})
  }
\]
is commutative; this means that for every $x\in M_0$, the sections
$\alpha_0^f(x)$ and $\alpha_0^g(x)$ of $\widetilde{M}$ coincide in
$D_+(f)\cap D_+(g)$, and consequently there exists a unique section
$\alpha_0(x)\in\Gamma(X,\widetilde{M})$ whose restriction to each
$D_+(f)$ is $\alpha_0^f(x)$. We have thus defined (without using the
hypothesis that $S$ is generated by $S_1$) a homomorphism of abelian groups
\[
\label{II.2.6.2.1}
  \alpha_0:M_0\longrightarrow\Gamma(X,\widetilde{M}).
\tag{2.6.2.1}
\]

Applying this result to the graded $S$-module $M(n)$ (for every
$n\in\mathbf{Z}$), we obtain for each $n\in\mathbf{Z}$ a homomorphism
of abelian groups
\[
\label{II.2.6.2.2}
  \alpha_n:M_n=(M(n))_0\longrightarrow\Gamma(X,\widetilde{M}(n))
\tag{2.6.2.2}
\]
(taking into account (\hyperref[II.2.5.15]{2.5.15})); whence a
functorial homomorphism (of degree $0$) of graded abelian groups
\[
\label{II.2.6.2.3}
  \alpha:M\longrightarrow\Gamma_*(\widetilde{M})
\tag{2.6.2.3}
\]
(also denoted by $\alpha_M$) which coincides with $\alpha_n$ on each $M_n$.

If in particular we take $M=S$, we verify immediately (taking into account
the definition of the multiplication in $\Gamma_*(\mathcal{O}_X)$
(\hyperref[0.5.4.6]{0, 5.4.6})) that
$\alpha:S\to\Gamma_*(\mathcal{O}_X)$ is a homomorphism of graded
rings, and that for every graded $S$-module $M$,
(\hyperref[II.2.6.2.3]{2.6.2.3}) is a di-homomorphism of graded
modules.
\end{env}

\begin{proposition}[2.6.3]
\label{II.2.6.3}
For every $f\in S_d$ ($d>0$), $D_+(f)$ is identical to the set of
$\mathfrak{p}\in X$ where the section $\alpha_d(f)$ of $\mathcal{O}_X(d)$
does not vanish (\hyperref[0.5.5.2]{0, 5.5.2}).
\end{proposition}

\begin{proof}
Since $X=\bigcup_{g\in S_1}D_+(g)$ by hypothesis, it suffices to verify that
for every $g\in S_1$, the set of $\mathfrak{p}\in D_+(g)$ where
$\alpha_d(f)$ does not vanish is identical to $D_+(fg)$. Now, the
restriction of $\alpha_d(f)$ to $D_+(g)$ is by definition the section
corresponding to the element $f/1$ of $(S(d))_{(g)}$; under the canonical
isomorphism
$(S(d))_{(g)}\xrightarrow{\sim}S_{(g)}$
(\hyperref[II.2.5.7]{2.5.7}), this section of $\mathcal{O}_X(d)$
over $D_+(g)$ is identified with the section of $\mathcal{O}_X$ over
$D_+(g)$ which corresponds to the element $f/g^d$ of $S_{(g)}$; to say that
this section vanishes at $\mathfrak{p}\in D_+(g)$ means that
$f/g^d\in\mathfrak{q}$, where $\mathfrak{q}$ is the prime ideal of $S_{(g)}$
corresponding to $\mathfrak{p}$ (\hyperref[II.2.3.6]{2.3.6}); by
definition this means that $f\in\mathfrak{p}$, whence the proposition.
\end{proof}

\begin{env}[2.6.4]
\label{II.2.6.4}
Now let $\mathcal{F}$ be an $\mathcal{O}_X$-module, and set
$M=\Gamma_*(\mathcal{F})$; by virtue of the existence of the homomorphism
of graded rings $\alpha:S\to\Gamma_*(\mathcal{O}_X)$, $M$ can be
regarded as a graded $S$-module. For every $f\in S_d$ ($d>0$), it
follows from (\hyperref[II.2.6.3]{2.6.3}) that the restriction to
$D_+(f)$ of the section $\alpha_d(f)$ of $\mathcal{O}_X(d)$ is invertible;
hence so is the restriction to $D_+(f)$ of the section $\alpha_d(f^n)$ of
$\mathcal{O}_X(nd)$, for every $n>0$. Let then
$z\in M_{nd}=\Gamma(X,\mathcal{F}(nd))$ ($n>0$); if there exists an integer
$k\geq0$ such that the restriction to $D_+(f)$ of $f^kz$, that is, the
\oldpage[II]{38}
section
$(z|D_+(f))(\alpha_d(f^k)|D_+(f))$ of $\mathcal{F}((n+k)d)$, is zero,
then, because of the preceding remark, we also have $z|D_+(f)=0$. This
shows that we define an $S_{(f)}$-homomorphism
$\beta_f:M_{(f)}\to\Gamma(D_+(f),\mathcal{F})$ by assigning to
the element $z/f^n$ the section
$(z|D_+(f))(\alpha_d(f^n)|D_+(f))^{-1}$ of $\mathcal{F}$ over
$D_+(f)$. Moreover, we verify immediately that the diagram
\[
\label{II.2.6.4.1}
  \xymatrix{
    M_{(f)}\ar[r]^{\beta_f}\ar[d]
    & \Gamma(D_+(f),\mathcal{F})\ar[d]\\
    M_{(fg)}\ar[r]_{\beta_{fg}}
    & \Gamma(D_+(fg),\mathcal{F})
  }
\tag{2.6.4.1}
\]
is commutative for $g\in S_e$ ($e>0$). If we recall that $M_{(f)}$
is canonically identified with $\Gamma(D_+(f),\widetilde{M})$ and that the
$D_+(f)$ form a basis of the topology of $X$
(\hyperref[II.2.3.4]{2.3.4}), we see that the $\beta_f$ come from a
unique canonical homomorphism of $\mathcal{O}_X$-modules
\[
\label{II.2.6.4.2}
  \beta:\widetilde{\Gamma_*(\mathcal{F})}
  \longrightarrow\mathcal{F}
\tag{2.6.4.2}
\]
(also denoted by $\beta_{\mathcal{F}}$), which is evidently functorial.
\end{env}

\begin{proposition}[2.6.5]
\label{II.2.6.5}
Let $M$ be a graded $S$-module and $\mathcal{F}$ an $\mathcal{O}_X$-module;
then the composite homomorphisms
\[
\label{II.2.6.5.1}
  \widetilde{M}\xrightarrow{\widetilde{\alpha}}
  \widetilde{\Gamma_*(\widetilde{M})}
  \xrightarrow{\beta}\widetilde{M}
\tag{2.6.5.1}
\]
\[
\label{II.2.6.5.2}
  \Gamma_*(\mathcal{F})\xrightarrow{\alpha}
  \Gamma_*(\widetilde{\Gamma_*(\mathcal{F})})
  \xrightarrow{\Gamma_*(\beta)}\Gamma_*(\mathcal{F})
\tag{2.6.5.2}
\]
are the identity isomorphisms.
\end{proposition}

\begin{proof}
The verification of (\hyperref[II.2.6.5.1]{2.6.5.1}) is local: in an
open subset $D_+(f)$, it follows immediately from the definitions and from
the fact that $\beta$, applied to quasi-coherent sheaves, is determined by
its action on the sections over $D_+(f)$ (\hyperref[I.1.3.8]{I, 1.3.8}).
The verification of (\hyperref[II.2.6.5.2]{2.6.5.2}) is done for each
degree separately: if we set $M=\Gamma_*(\mathcal{F})$, we have
$M_n=\Gamma(X,\mathcal{F}(n))$ and
$(\Gamma_*(\widetilde{M}))_n=\Gamma(X,\widetilde{M}(n))
=\Gamma(X,\widetilde{M(n)})$. Now, if $f\in S_1$ and $z\in M_n$,
$\alpha_n^f(z)$ is the element $z/1$ of $(M(n))_{(f)}$, equal to
$(f/1)^n(z/f^n)$; to it there corresponds under $\beta_f$ the section
\[
  ((\alpha_1(f))^n|D_+(f))
  ((z|D_+(f))((\alpha_1(f))^n|D_+(f))^{-1})
\]
over $D_+(f)$, that is, the restriction of $z$ to $D_+(f)$, which
verifies (\hyperref[II.2.6.5.2]{2.6.5.2}).
\end{proof}


\subsection{Finiteness conditions.}
\label{subsection:II.2.7}

\begin{proposition}[2.7.1]
\label{II.2.7.1}
\begin{enumerate}
  \item[\rm(i)] If $S$ is a Noetherian graded ring,
    $X=\operatorname{Proj}(S)$ is a Noetherian scheme.
  \item[\rm(ii)] If $S$ is a graded $A$-algebra of finite type,
    $X=\operatorname{Proj}(S)$ is a scheme of finite type over
    $Y=\operatorname{Spec}(A)$.
\end{enumerate}
\end{proposition}

\oldpage[II]{39}
\begin{proof}
\begin{enumerate}
  \item[\rm(i)] If $S$ is Noetherian, the ideal $S_+$ admits a finite system of
    homogeneous generators $f_i$ ($1\leq i\leq p$), so
    (\hyperref[II.2.3.14]{2.3.14}) the underlying space $X$ is the union of
    the $D_+(f_i)=\operatorname{Spec}(S_{(f_i)})$, and everything comes down to
    showing that each of the $S_{(f_i)}$ is Noetherian, which follows from
    (\hyperref[II.2.2.6]{2.2.6}).
  \item[\rm(ii)] The hypothesis implies that $S_0$ is an $A$-algebra of finite
    type and $S$ an $S_0$-algebra of finite type, so $S_+$ is an ideal of
    finite type (\hyperref[II.2.1.4]{2.1.4}). We are then reduced, as in
    (i), to proving that for $f\in S_d$, $S_{(f)}$ is an $A$-algebra of finite
    type. By virtue of (\hyperref[II.2.2.5]{2.2.5}), it suffices to show
    that $S^{(d)}$ is an $A$-algebra of finite type, which follows from
    (\hyperref[II.2.1.6]{2.1.6}).
\end{enumerate}
\end{proof}

\begin{env}[2.7.2]
\label{II.2.7.2}
In what follows we will consider the following finiteness conditions for
a graded $S$-module $M$:
\begin{enumerate}
  \item[\rm(TF)] There exists an integer $n$ such that the submodule
    $\bigoplus_{k\geq n}M_k$ is an $S$-module of finite type.
  \item[\rm(TN)] There exists an integer $n$ such that $M_k=0$ for $k\geq n$.
\end{enumerate}

If $M$ satisfies (TN), we have $M_{(f)}=0$ for every homogeneous $f$ in $S_+$,
hence $\widetilde{M}=0$.

Let $M$, $N$ be two graded $S$-modules; we say that a homomorphism
$u:M\to N$ of degree $0$ is (TN)-\emph{injective} (resp.
(TN)-\emph{surjective}, (TN)-\emph{bijective}) if there exists an integer $n$
such that $u_k:M_k\to N_k$ is injective (resp. surjective, bijective) for
$k\geq n$. To say that $u$ is (TN)-injective (resp. (TN)-surjective) therefore
amounts to saying that $\operatorname{Ker}u$ (resp. $\operatorname{Coker}u$)
satisfies (TN). By virtue of (\hyperref[II.2.5.4]{2.5.4}), if $u$ is
(TN)-injective (resp. (TN)-surjective, (TN)-bijective), $\widetilde{u}$ is
injective (resp. surjective, bijective); when $u$ is (TN)-bijective, we also
say that $u$ is a (TN)-\emph{isomorphism}.
\end{env}

\begin{proposition}[2.7.3]
\label{II.2.7.3}
Let $S$ be a graded ring such that the ideal $S_+$ is of finite type, and $M$ a
graded $S$-module.
\begin{enumerate}
  \item[\rm(i)] If $M$ satisfies condition (TF), the $\mathcal{O}_X$-module
    $\widetilde{M}$ is of finite type.
  \item[\rm(ii)] Suppose that $M$ satisfies (TF); in order that
    $\widetilde{M}=0$, it is necessary and sufficient that $M$ satisfy (TN).
\end{enumerate}
\end{proposition}

\begin{proof}
We have just seen that condition (TN) implies $\widetilde{M}=0$. If $M$
satisfies (TF), the graded submodule $M'=\bigoplus_{k\geq n}M_k$, which is of
finite type by hypothesis, is such that $M/M'$ satisfies (TN); consequently
$\widetilde{M/M'}=0$, and the exactness of the functor
$M\mapsto\widetilde{M}$ (\hyperref[II.2.5.4]{2.5.4}) implies
$\widetilde{M}=\widetilde{M'}$; to prove that $\widetilde{M}$ is of finite
type, we are therefore reduced to the case where $M$ is of finite type. The
question being local, it suffices to prove that $M_{(f)}$ is an
$S_{(f)}$-module of finite type (\hyperref[I.1.3.9]{I, 1.3.9}); but
$M^{(d)}$ is an $S^{(d)}$-module of finite type
(\hyperref[II.2.1.6]{2.1.6, (iii)}) and our assertion then follows from
(\hyperref[II.2.2.5]{2.2.5}).

Suppose now that $M$ satisfies (TF) and that $\widetilde{M}=0$; then, with
the same notation as above, we have $\widetilde{M'}=0$, and condition (TN)
for $M'$ is equivalent to condition (TN) for $M$; so, to prove that
$\widetilde{M}=0$ implies that $M$ satisfies (TN), we may again restrict to
the case where $M$ is generated by finitely many homogeneous elements $x_i$
($1\leq i\leq p$); on the other hand, let $(f_j)_{1\leq j\leq q}$ be a system
of homogeneous generators of the ideal $S_+$. By hypothesis we have
$M_{(f_j)}=0$ for every $j$, so there exists an integer $n$ such that
$f_j^nx_i=0$ for all $i$ and $j$. Let $n_j$ be the degree of $f_j$, and let
$m$ be the largest value of $\sum_jr_jn_j$ for the systems of integers
$(r_j)$, finite in number, such that $\sum_jr_j\leq nq$; it is then clear

\oldpage[II]{40}
that if $k>m$, we have $S_kx_i=0$ for every $i$; if $h$ is the largest of the
degrees of the $x_i$, we conclude that $M_k=0$ for $k>h+m$, which completes
the proof.
\end{proof}

\begin{corollary}[2.7.4]
\label{II.2.7.4}
Let $S$ be a graded ring such that the ideal $S_+$ is of finite type; in order
that $X=\operatorname{Proj}(S)=\varnothing$, it is necessary and sufficient
that there exist $n$ such that $S_k=0$ for $k\geq n$.
\end{corollary}

\begin{proof}
Indeed, the condition $X=\varnothing$ is equivalent to
$\mathcal{O}_X=\widetilde{S}=0$, and $S$ is a monogenic $S$-module.
\end{proof}

\begin{theorem}[2.7.5]
\label{II.2.7.5}
Suppose that the ideal $S_+$ is generated by finitely many homogeneous
elements of degree $1$; let $X=\operatorname{Proj}(S)$. Then, for every
quasi-coherent $\mathcal{O}_X$-module $\mathcal{F}$, the canonical
homomorphism $\beta:\widetilde{\Gamma_*(\mathcal{F})}
\to\mathcal{F}$ (\hyperref[II.2.6.4]{2.6.4}) is an isomorphism.
\end{theorem}

\begin{proof}
Indeed, if $S_+$ is generated by finitely many elements $f_i\in S_1$, $X$ is
the union of the subspaces $\operatorname{Spec}(S_{(f_i)})$
(\hyperref[II.2.3.6]{2.3.6}), which are quasi-compact, and is therefore
quasi-compact; moreover $X$ is a scheme
(\hyperref[II.2.4.2]{2.4.2}); by
(\hyperref[I.9.3.2]{I, 9.3.2}),
(\hyperref[II.2.5.14.2]{2.5.14.2}) and
(\hyperref[II.2.6.3]{2.6.3}) we have, for every $f\in S_d$ ($d>0$), a
canonical isomorphism
$(\Gamma_*(\mathcal{F}))_{(\alpha_d(f))}\xrightarrow{\sim}
\Gamma(D_+(f),\mathcal{F})$; besides, by definition,
$(\Gamma_*(\mathcal{F}))_{(\alpha_d(f))}$ (where $\Gamma_*(\mathcal{F})$ is
regarded as a $\Gamma_*(\mathcal{O}_X)$-module) is none other than
$(\Gamma_*(\mathcal{F}))_{(f)}$ (where $\Gamma_*(\mathcal{F})$ is regarded
as an $S$-module); if we refer to the definition
(\hyperref[I.9.3.1]{I, 9.3.1}) of the preceding canonical isomorphism, we
see that it coincides with the homomorphism $\beta_f$, whence the theorem.
\end{proof}

\begin{remark}[2.7.6]
\label{II.2.7.6}
If the graded ring $S$ is assumed to be \emph{Noetherian}, the condition of
(\hyperref[II.2.7.5]{2.7.5}) is satisfied \emph{ipso facto} as soon as the
ideal $S_+$ is assumed to be generated by the set $S_1$ of homogeneous elements
of degree $1$.
\end{remark}

\begin{corollary}[2.7.7]
\label{II.2.7.7}
Under the hypotheses of (\hyperref[II.2.7.5]{2.7.5}), every
quasi-coherent $\mathcal{O}_X$-module $\mathcal{F}$ is isomorphic to an
$\mathcal{O}_X$-module of the form $\widetilde{M}$, where $M$ is a graded
$S$-module.
\end{corollary}

\begin{corollary}[2.7.8]
\label{II.2.7.8}
Under the hypotheses of (\hyperref[II.2.7.5]{2.7.5}), every
quasi-coherent $\mathcal{O}_X$-module $\mathcal{F}$ of finite type is
isomorphic to an $\mathcal{O}_X$-module of the form $\widetilde{N}$, where $N$
is a graded $S$-module of finite type.
\end{corollary}

\begin{proof}
We may suppose $\mathcal{F}=\widetilde{M}$, where $M$ is a graded $S$-module
(\hyperref[II.2.7.7]{2.7.7}). Let $(f_\lambda)_{\lambda\in L}$ be a
system of homogeneous generators of $M$; for every finite subset $H$ of $L$,
let $M_H$ be the graded submodule of $M$ generated by the $f_\lambda$ such
that $\lambda\in H$; it is clear that $M$ is the inductive limit of its
submodules $M_H$, so $\mathcal{F}$ is the inductive limit of its
sub-$\mathcal{O}_X$-modules $\widetilde{M_H}$ (\hyperref[II.2.5.4]{2.5.4}).
But since $\mathcal{F}$ is of finite type and the underlying space $X$ is
quasi-compact, it follows from (\hyperref[0.5.2.3]{0, 5.2.3}) that
$\mathcal{F}=\widetilde{M_H}$ for some finite subset $H$ of $L$.
\end{proof}

\begin{corollary}[2.7.9]
\label{II.2.7.9}
Under the hypotheses of (\hyperref[II.2.7.5]{2.7.5}), let $\mathcal{F}$ be
a quasi-coherent $\mathcal{O}_X$-module of finite type. Then there exists an
integer $n_0$ such that, for every $n\geq n_0$, $\mathcal{F}(n)$ is isomorphic
to a quotient of an $\mathcal{O}_X$-module of the form $\mathcal{O}_X^k$ ($k>0$
depending on $n$), and consequently is generated by finitely many of its
sections over $X$ (\hyperref[0.5.1.1]{0, 5.1.1}).
\end{corollary}

\begin{proof}
By virtue of (\hyperref[II.2.7.8]{2.7.8}), we may suppose that
$\mathcal{F}=\widetilde{M}$, where $M$ is a quotient of a finite direct sum of
$S$-modules of the form $S(m_i)$; by virtue of
(\hyperref[II.2.5.4]{2.5.4}) we may therefore restrict to the case where
$M=S(m)$, so that $\mathcal{F}(n)=\widetilde{S(m+n)}
=\mathcal{O}_X(m+n)$ (\hyperref[II.2.5.15]{2.5.15}). It will therefore
suffice to prove

\oldpage[II]{41}
\begin{lemma}[2.7.9.1]
\label{II.2.7.9.1}
Under the hypotheses of (\hyperref[II.2.7.5]{2.7.5}), for every $n\geq0$,
there exist an integer $k$ (depending on $n$) and a surjective homomorphism
$\mathcal{O}_X^k\to\mathcal{O}_X(n)$.
\end{lemma}

Indeed, it suffices (\hyperref[II.2.7.2]{2.7.2}) to show that, for a
suitable $k$, there is a (TN)-surjective homomorphism $u$ \emph{of degree $0$}
from the product graded $S$-module $S^k$ to $S(n)$. Now, we have
$(S(n))_0=S_n$, and by hypothesis $S_h=S_1^h$ for every $h>0$, so
$SS_n=S_n+S_{n+1}+\cdots$. Since $S_n$ is an $S_0$-module of finite type
(\hyperref[II.2.1.5]{2.1.5} and
\hyperref[II.2.1.6]{2.1.6, (i)}), consider a system of generators
$(a_i)_{1\leq i\leq k}$ of this module; the homomorphism $u$ will send the
$i$-th element of the canonical basis of $S^k$ to $a_i$ ($1\leq i\leq k$);
since $\operatorname{Coker}u$ is then identified with
$(S(n))_{-n}+\cdots+(S(n))_{-1}$, $u$ indeed meets the requirement.
\end{proof}

\begin{corollary}[2.7.10]
\label{II.2.7.10}
Under the hypotheses of (\hyperref[II.2.7.5]{2.7.5}), let $\mathcal{F}$ be
a quasi-coherent $\mathcal{O}_X$-module of finite type. There exists an
integer $n_0$ such that, for every $n\geq n_0$, $\mathcal{F}$ is isomorphic to
a quotient of an $\mathcal{O}_X$-module of the form $(\mathcal{O}_X(-n))^k$
($k$ depending on $n$).
\end{corollary}

\begin{proposition}[2.7.11]
\label{II.2.7.11}
Suppose that the hypotheses of
(\hyperref[II.2.7.5]{2.7.5}) are satisfied, and let $M$ be a graded
$S$-module. Then:
\begin{enumerate}
  \item[\rm(i)] The canonical homomorphism
    $\widetilde{\alpha}:\widetilde{M}\to
    \widetilde{\Gamma_*(\widetilde{M})}$ is an isomorphism.
  \item[\rm(ii)] Let $\mathcal{G}$ be a quasi-coherent
    sub-$\mathcal{O}_X$-module of $\widetilde{M}$, and let $N$ be the graded
    sub-$S$-module of $M$ that is the inverse image of $\Gamma_*(\mathcal{G})$
    under $\alpha$. Then we have $\widetilde{N}=\mathcal{G}$ ($\widetilde{N}$
    being identified, by virtue of (\hyperref[II.2.5.4]{2.5.4}), with a
    sub-$\mathcal{O}_X$-module of $\widetilde{M}$).
\end{enumerate}
\end{proposition}

\begin{proof}
Since $\beta:\widetilde{\Gamma_*(\widetilde{M})}
\to\widetilde{M}$ is an isomorphism by virtue of
(\hyperref[II.2.7.5]{2.7.5}), $\widetilde{\alpha}$ is the inverse
isomorphism by virtue of (\hyperref[II.2.6.5.1]{2.6.5.1}), whence (i). Let
$P$ be the graded submodule $\alpha(M)$ of $\Gamma_*(\widetilde{M})$; since
$M\mapsto\widetilde{M}$ is an exact functor
(\hyperref[II.2.5.4]{2.5.4}), the image of $\widetilde{M}$ under
$\widetilde{\alpha}$ is equal to $\widetilde{P}$, so, by virtue of (i),
$\widetilde{P}=\widetilde{\Gamma_*(\widetilde{M})}$. Set
$Q=\Gamma_*(\mathcal{G})\cap P$; by virtue of the foregoing and of
(\hyperref[II.2.5.4]{2.5.4}), we have
$\widetilde{Q}=\widetilde{\Gamma_*(\mathcal{G})}$, so the restriction of
$\beta$ to $\widetilde{Q}$ is an \emph{isomorphism} of this
$\mathcal{O}_X$-module onto $\mathcal{G}$ by
(\hyperref[II.2.7.5]{2.7.5}). But by the definition of $N$ and
(\hyperref[II.2.5.4]{2.5.4}), the restriction of the isomorphism
$\widetilde{\alpha}$ to $\widetilde{N}$ is an isomorphism of $\widetilde{N}$
onto $\widetilde{Q}$, whence the conclusion by
(\hyperref[II.2.6.5.1]{2.6.5.1}).
\end{proof}

\subsection{Functorial behaviour.}
\label{subsection:II.2.8}

\begin{env}[2.8.1]
\label{II.2.8.1}
Let $S$, $S'$ be two rings graded in positive degrees, and
$\varphi:S'\to S$ a homomorphism of graded rings. We denote by
$G(\varphi)$ the open subset of $X=\operatorname{Proj}(S)$ complementary to
$V_+(\varphi(S'_+))$, or, what amounts to the same, the union of the
$D_+(\varphi(f'))$, where $f'$ runs through the set of homogeneous elements of
$S'_+$. The restriction to $G(\varphi)$ of the continuous map
${}^a\varphi$ from $\operatorname{Spec}(S)$ to $\operatorname{Spec}(S')$
(\hyperref[I.1.2.1]{I, 1.2.1}) is therefore a continuous map from
$G(\varphi)$ to $\operatorname{Proj}(S')$, which we again denote by
${}^a\varphi$ by abuse of language. If $f'\in S'_+$ is homogeneous, we have
\[
\label{II.2.8.1.1}
  {}^a\varphi^{-1}(D_+(f'))=D_+(\varphi(f'))
\tag{2.8.1.1}
\]
taking into account the fact that ${}^a\varphi$ maps $G(\varphi)$ into
$\operatorname{Proj}(S')$, and (\hyperref[I.1.2.2.2]{I, 1.2.2.2}). On the
other hand, the homomorphism $\varphi$ canonically defines (with the same
notation) a homomorphism of graded rings $S'_{f'}\to S_f$, whence, by
restriction to the elements of degree $0$,

\oldpage[II]{42}
a homomorphism $S'_{(f')}\to S_{(f)}$, which we denote by
$\varphi_{(f)}$; to it there corresponds
(\hyperref[I.1.6.1]{I, 1.6.1}) a morphism
$({}^a\varphi_{(f)},\widetilde{\varphi}_{(f)}):
\operatorname{Spec}(S_{(f)})\to\operatorname{Spec}(S'_{(f')})$ of affine
schemes. If we canonically identify $\operatorname{Spec}(S_{(f)})$ with the
scheme induced by $\operatorname{Proj}(S)$ on $D_+(f)$
(\hyperref[II.2.3.6]{2.3.6}), we have defined a morphism
$\Phi_f:D_+(f)\to D_+(f')$, and ${}^a\varphi_{(f)}$ is identified with the
restriction of ${}^a\varphi$ to $D_+(f)$. On the other hand, it is immediate
that if $g'$ is a second homogeneous element of $S'_+$ and $g=\varphi(g')$,
the diagram
\[
  \xymatrix{
    D_+(f)\ar[r]^{\Phi_f} & D_+(f')\\
    D_+(fg)\ar[u]\ar[r]_{\Phi_{fg}} & D_+(f'g')\ar[u]
  }
\]
is commutative, because of the commutativity of the diagram
\[
  \xymatrix{
    S'_{(f')}\ar[r]^{\varphi_{(f)}}\ar[d]_{\omega_{f'g',f'}}
    & S_{(f)}\ar[d]^{\omega_{fg,f}}\\
    S'_{(f'g')}\ar[r]_{\varphi_{(fg)}} & S_{(fg)}
  }
\]
Taking into account the definition of $G(\varphi)$ and
(\hyperref[II.2.3.3.2]{2.3.3.2}), we therefore see that:
\end{env}

\begin{proposition}[2.8.2]
\label{II.2.8.2}
Given a homomorphism of graded rings $\varphi:S'\to S$, there exists one and
only one morphism $({}^a\varphi,\widetilde{\varphi})$ from the induced
prescheme $G(\varphi)$ to $\operatorname{Proj}(S')$ (said to be
\emph{associated to $\varphi$} and denoted $\operatorname{Proj}(\varphi)$) such
that, for every homogeneous element $f'\in S'_+$, the restriction of this
morphism to $D_+(\varphi(f'))$ coincides with the morphism associated to the
homomorphism $S'_{(f')}\to S_{(\varphi(f'))}$ corresponding to $\varphi$.
\end{proposition}

\begin{proof}
We note that, with the preceding notation, if $f'\in S'_d$, the diagram
\[
\label{II.2.8.2.1}
  \xymatrix{
    S'_{(f')}\ar[r]^{\varphi_{(f)}}\ar[d]_{\sim}
    & S_{(f)}\ar[d]^{\sim}\\
    {S'}^{(d)}/(f'-1){S'}^{(d)}\ar[r]
    & S^{(d)}/(f-1)S^{(d)}
  }
\tag{2.8.2.1}
\]
is commutative (the vertical arrows being the isomorphisms
(\hyperref[II.2.2.5]{2.2.5})).
\end{proof}

\begin{corollary}[2.8.3]
\label{II.2.8.3}
\begin{enumerate}
  \item[\rm(i)] The morphism $\operatorname{Proj}(\varphi)$ is affine.
  \item[\rm(ii)] If $\operatorname{Ker}(\varphi)$ is nilpotent (and in
    particular if $\varphi$ is injective), the morphism
    $\operatorname{Proj}(\varphi)$ is dominant.
\end{enumerate}
\end{corollary}

\begin{proof}
(i) is an immediate consequence of (\hyperref[II.2.8.2]{2.8.2}) and
(\hyperref[II.2.8.1.1]{2.8.1.1}). On the other hand, if
$\operatorname{Ker}(\varphi)$ is nilpotent, then for every homogeneous $f'$ in
$S'_+$ we verify immediately that $\operatorname{Ker}(\varphi_f)$ (with
$f=\varphi(f')$) is nilpotent, hence so is
$\operatorname{Ker}(\varphi_{(f)})$; the conclusion therefore follows from
(\hyperref[II.2.8.2]{2.8.2}) and
(\hyperref[I.1.2.7]{I, 1.2.7}).
\end{proof}

\oldpage[II]{43}
We note that there are in general morphisms from $\operatorname{Proj}(S)$
to $\operatorname{Proj}(S')$ that are not affine, and therefore do not come
from homomorphisms of graded rings $S'\to S$; an example is the structural
morphism $\operatorname{Proj}(S)\to\operatorname{Spec}(A)$ when $A$ is a field
($\operatorname{Spec}(A)$ being identified with
$\operatorname{Proj}(A[T])$ (cf. (\hyperref[II.3.1.7]{3.1.7}))); indeed,
this follows from (\hyperref[I.2.3.2]{I, 2.3.2}).

\begin{env}[2.8.4]
\label{II.2.8.4}
Let $S''$ be a third ring graded in positive degrees,
$\varphi':S''\to S'$ a homomorphism of graded rings, and set
$\varphi''=\varphi\circ\varphi'$. Because of
(\hyperref[II.2.8.1.1]{2.8.1.1}) and the formula
${}^a\varphi''=({}^a\varphi')\circ({}^a\varphi)$, we verify immediately that
$G(\varphi'')\subset G(\varphi)$ and that, if $\Phi$, $\Phi'$ and $\Phi''$
are the morphisms associated to $\varphi$, $\varphi'$ and $\varphi''$, we have
$\Phi''=\Phi'\circ(\Phi|G(\varphi''))$.
\end{env}

\begin{env}[2.8.5]
\label{II.2.8.5}
Suppose that $S$ (resp. $S'$) is a graded $A$-algebra (resp. a graded
$A'$-algebra), and let $\psi:A'\to A$ be a ring homomorphism such that the
diagram
\[
\label{II.2.8.5.1}
  \xymatrix{
    A'\ar[r]^{\psi}\ar[d] & A\ar[d]\\
    S'\ar[r]_{\varphi} & S
  }
\tag{2.8.5.1}
\]
is commutative. We can then regard $G(\varphi)$ and
$\operatorname{Proj}(S')$ as schemes over $\operatorname{Spec}(A)$ and
$\operatorname{Spec}(A')$ respectively; if $\Phi$ and $\Psi$ are the morphisms
associated respectively to $\varphi$ and $\psi$, the diagram
\[
  \xymatrix{
    G(\varphi)\ar[r]^{\Phi}\ar[d] & \operatorname{Proj}(S')\ar[d]\\
    \operatorname{Spec}(A)\ar[r]_{\Psi} & \operatorname{Spec}(A')
  }
\]
is then commutative: indeed, it suffices to prove this for the restriction of
$\Phi$ to $D_+(f)$, where $f=\varphi(f')$, $f'$ homogeneous in $S'_+$; and
then it follows from the commutativity of the diagram
\[
  \xymatrix{
    A'\ar[r]^{\psi}\ar[d] & A\ar[d]\\
    S'_{(f')}\ar[r]_{\varphi_{(f)}} & S_{(f)}
  }
\]
\end{env}

\begin{env}[2.8.6]
\label{II.2.8.6}
Now let $M$ be a graded $S$-module, and consider the $S'$-module
$M_{[\varphi]}$, which is evidently graded. Let $f'$ be homogeneous in
$S'_+$, and let $f=\varphi(f')$; we know
(\hyperref[0.1.5.2]{0, 1.5.2}) that there is a canonical isomorphism
$(M_{[\varphi]})_{f'}\xrightarrow{\sim}(M_f)_{[\varphi_f]}$, and it is
immediate that this isomorphism preserves degrees, so there is a canonical
isomorphism
$(M_{[\varphi]})_{(f')}\xrightarrow{\sim}
(M_{(f)})_{[\varphi_{(f)}]}$. To this isomorphism there canonically
corresponds an isomorphism of sheaves
$\widetilde{M_{[\varphi]}}|D_+(f')
\xrightarrow{\sim}(\Phi_f)_*(\widetilde{M}|D_+(f))$
(\hyperref[II.2.5.2]{2.5.2} and
\hyperref[I.1.6.3]{I, 1.6.3}). Moreover,

\oldpage[II]{44}
if $g'$ is a second homogeneous element of $S'_+$ and $g=\varphi(g')$, the
diagram
\[
  \xymatrix{
    (M_{[\varphi]})_{(f')}\ar[r]^{\sim}\ar[d]
    & (M_{(f)})_{[\varphi_{(f)}]}\ar[d]\\
    (M_{[\varphi]})_{(f'g')}\ar[r]^{\sim}
    & (M_{(fg)})_{[\varphi_{(fg)}]}
  }
\]
is commutative, from which we immediately conclude that the isomorphism
\[
  \widetilde{M_{[\varphi]}}|D_+(f'g')
  \xrightarrow{\sim}(\Phi_{fg})_*(\widetilde{M}|D_+(fg))
\]
is the restriction to $D_+(f'g')$ of the isomorphism
$\widetilde{M_{[\varphi]}}|D_+(f')
\xrightarrow{\sim}(\Phi_f)_*(\widetilde{M}|D_+(f))$. Since $\Phi_f$ is the
restriction to $D_+(f)$ of the morphism $\Phi$, we therefore see, taking into
account (\hyperref[II.2.8.1.1]{2.8.1.1}), and setting
$X'=\operatorname{Proj}(S')$:
\end{env}

\begin{proposition}[2.8.7]
\label{II.2.8.7}
There exists a canonical functorial isomorphism of the
$\mathcal{O}_{X'}$-module $\widetilde{M_{[\varphi]}}$ onto the
$\mathcal{O}_{X'}$-module $\Phi_*(\widetilde{M}|G(\varphi))$.
\end{proposition}

From this we immediately deduce a canonical functorial map from the set of
$\varphi$-morphisms $M'\to M$ of a graded $S'$-module into the graded
$S$-module $M$ to the set of $\Phi$-morphisms
$\widetilde{M'}\to\widetilde{M}|G(\varphi)$. With the notation of
(\hyperref[II.2.8.4]{2.8.4}), if $M''$ is a graded $S''$-module, to the
composite of a $\varphi$-morphism $M'\to M$ and a $\varphi'$-morphism
$M''\to M'$ there canonically corresponds the composite of
$\widetilde{M'}|G(\varphi')\to\widetilde{M}|G(\varphi'')$ and
$\widetilde{M''}\to\widetilde{M'}|G(\varphi')$.

\begin{proposition}[2.8.8]
\label{II.2.8.8}
Under the hypotheses of (\hyperref[II.2.8.1]{2.8.1}), let $M'$ be a graded
$S'$-module. There exists a canonical functorial homomorphism $\nu$ from the
$(\mathcal{O}_X|G(\varphi))$-module $\Phi^*(\widetilde{M'})$ to the
$(\mathcal{O}_X|G(\varphi))$-module
$\widetilde{M'\otimes_{S'}S}|G(\varphi)$. If the ideal $S'_+$
is generated by $S'_1$, $\nu$ is an isomorphism.
\end{proposition}

\begin{proof}
Indeed, for $f'\in S'_d$ ($d>0$), we define a canonical functorial
homomorphism of $S_{(f)}$-modules (where $f=\varphi(f')$)
\[
\label{II.2.8.8.1}
  \nu_f:M'_{(f')}\otimes_{S'_{(f')}}S_{(f)}
  \longrightarrow(M'\otimes_{S'}S)_{(f)}
\tag{2.8.8.1}
\]
by composing the homomorphism
$M'_{(f')}\otimes_{S'_{(f')}}S_{(f)}\to
M'_{f'}\otimes_{S'_{f'}}S_f$ and the canonical isomorphism
$M'_{f'}\otimes_{S'_{f'}}S_f\xrightarrow{\sim}(M'\otimes_{S'}S)_f$
(\hyperref[0.1.5.4]{0, 1.5.4}), and noting that the latter preserves
degrees. We verify immediately that the $\nu_f$ are compatible with the
restriction operators from $D_+(f)$ to $D_+(fg)$ (for a second $g'\in S'_+$
and $g=\varphi(g')$), whence the definition of the homomorphism
\[
  \nu:\Phi^*(\widetilde{M'})
  \longrightarrow\widetilde{M'\otimes_{S'}S}|G(\varphi)
\]
taking into account (\hyperref[I.1.6.5]{I, 1.6.5}). To prove the second
assertion, it suffices to show that $\nu_f$ is an isomorphism for every
$f'\in S'_1$, since $G(\varphi)$ is then the union of the
$D_+(\varphi(f'))$. We first define a $\mathbf{Z}$-bilinear map
$M'_m\times S_n\to M'_{(f')}\otimes_{S'_{(f')}}S_{(f)}$ by sending
$(x',s)$ to the element
$(x'/{f'}^m)\otimes(s/f^n)$ (with the convention that $x'/{f'}^m$ is
${f'}^{-m}x'/1$

\oldpage[II]{45}
when $m<0$). We see, as in the proof of
(\hyperref[II.2.5.13]{2.5.13}), that this map gives rise to a
di-homomorphism of modules
\[
  \eta_f:M'\otimes_{S'}S
  \longrightarrow M'_{(f')}\otimes_{S'_{(f')}}S_{(f)}.
\]
Moreover, if, for $r>0$, we have
$f^r\sum_i(x'_i\otimes s_i)=0$, this can also be written
$\sum_i({f'}^rx'_i\otimes s_i)=0$, whence, by
(\hyperref[0.1.5.4]{0, 1.5.4}),
$\sum_i({f'}^rx'_i/{f'}^{m_i+r})\otimes(s_i/f^{n_i})=0$, that is,
$\eta_f(\sum_i x'_i\otimes s_i)=0$, which proves that $\eta_f$ factors
as
\[
  M'\otimes_{S'}S\longrightarrow(M'\otimes_{S'}S)_f
  \xrightarrow{\eta'_f}M'_{(f')}\otimes_{S'_{(f')}}S_{(f)}~;
\]
finally, we verify that $\eta'_f$ and $\nu_f$ are isomorphisms inverse to each
other.

In particular, it follows from (\hyperref[II.2.1.2.1]{2.1.2.1}) that we
have a canonical homomorphism
\[
\label{II.2.8.8.2}
  \Phi^*(\mathcal{O}_{X'}(n))\xrightarrow{\sim}
  \mathcal{O}_X(n)|G(\varphi)
\tag{2.8.8.2}
\]
for every $n\in\mathbf{Z}$.
\end{proof}

\begin{env}[2.8.9]
\label{II.2.8.9}
Let $A$, $A'$ be two rings, and $\psi:A'\to A$ a ring homomorphism, defining
a morphism
$\Psi:\operatorname{Spec}(A)\to\operatorname{Spec}(A')$. Let $S'$ be an
$A'$-algebra graded in positive degrees, and set $S=S'\otimes_{A'}A$, which
is, in an evident way, an $A$-algebra graded by the $S'_n\otimes_{A'}A$;
the map $\varphi:s'\mapsto s'\otimes1$ is then a homomorphism of graded rings
making the diagram (\hyperref[II.2.8.5.1]{2.8.5.1}) commutative. Since here
$S_+$ is the $A$-module generated by $\varphi(S'_+)$, we have
$G(\varphi)=\operatorname{Proj}(S)=X$; whence, setting
$X'=\operatorname{Proj}(S')$, a commutative diagram
\[
\label{II.2.8.9.1}
  \xymatrix{
    X\ar[r]^{\Phi}\ar[d]_p & X'\ar[d]\\
    Y\ar[r]_{\Psi} & Y'
  }
\tag{2.8.9.1}
\]

Now let $M'$ be a graded $S'$-module, and set
$M=M'\otimes_{A'}A=M'\otimes_{S'}S$. Under these conditions:
\end{env}

\begin{proposition}[2.8.10]
\label{II.2.8.10}
The diagram (\hyperref[II.2.8.9.1]{2.8.9.1}) identifies the scheme $X$ with
the product $X'\times_{Y'}Y$; moreover, the canonical homomorphism
$\nu:\Phi^*(\widetilde{M'})\to\widetilde{M}$
(\hyperref[II.2.8.8]{2.8.8}) is an isomorphism.
\end{proposition}

\begin{proof}
The first assertion will be established if we prove that, for every
homogeneous $f'$ in $S'_+$, setting $f=\varphi(f')$, the restrictions of
$\Phi$ and $p$ to $D_+(f)$ identify this scheme with the product
$D_+(f')\times_{Y'}Y$ (\hyperref[I.3.2.6.2]{I, 3.2.6.2}); in other words,
it suffices (\hyperref[I.3.2.2]{I, 3.2.2}) to prove that $S_{(f)}$ is
canonically identified with $S'_{(f')}\otimes_{A'}A$, which is immediate by
virtue of the existence of the canonical degree-preserving isomorphism
$S_f\xrightarrow{\sim}S'_{f'}\otimes_{A'}A$
(\hyperref[0.1.5.4]{0, 1.5.4}). The second assertion follows from the fact
that $M'_{(f')}\otimes_{S'_{(f')}}S_{(f)}$ is identified, by the foregoing,
with $M'_{(f')}\otimes_{A'}A$, and the latter with $M_{(f)}$, since $M_f$ is
canonically identified with $M'_{f'}\otimes_{A'}A$ by a degree-preserving
isomorphism.
\end{proof}

\begin{corollary}[2.8.11]
\label{II.2.8.11}
For every integer $n\in\mathbf{Z}$, $\widetilde{M(n)}$ is identified with
$\Phi^*(\widetilde{M'(n)})=\widetilde{M'(n)}\otimes_{Y'}\mathcal{O}_Y$;
in particular
$\mathcal{O}_X(n)=\Phi^*(\mathcal{O}_{X'}(n))
=\mathcal{O}_{X'}(n)\otimes_{Y'}\mathcal{O}_Y$.
\end{corollary}

\begin{proof}
This follows from (\hyperref[II.2.8.10]{2.8.10}) and
(\hyperref[II.2.5.15]{2.5.15}).
\end{proof}

\oldpage[II]{46}

\begin{env}[2.8.12]
\label{II.2.8.12}
Under the hypotheses of (\hyperref[II.2.8.9]{2.8.9}), for
$f'\in S'_d$ ($d>0$) and $f=\varphi(f')$, the diagram
\[
\label{II.2.8.12.1}
  \xymatrix{
    M'_{(f')}\ar[r]^{\sim}\ar[d]
    & M'^{(d)}/(f'-1)M'^{(d)}\ar[d]\\
    M_{(f)}\ar[r]^{\sim}
    & M^{(d)}/(f-1)M^{(d)}
  }
\tag{2.8.12.1}
\]
(cf. (\hyperref[II.2.2.5]{2.2.5})) is commutative.
\end{env}

\begin{env}[2.8.13]
\label{II.2.8.13}
We keep the notation and hypotheses of (\hyperref[II.2.8.9]{2.8.9}), and
let $\mathcal{F}'$ be an $\mathcal{O}_{X'}$-module; if we set
$\mathcal{F}=\Phi^*(\mathcal{F}')$, we have, for every $n\in\mathbf{Z}$,
$\mathcal{F}(n)=\Phi^*(\mathcal{F}'(n))$ by virtue of
(\hyperref[II.2.8.11]{2.8.11}) and
(\hyperref[0.4.3.3]{0, 4.3.3}). Consequently
(\hyperref[0.3.7.1]{0, 3.7.1}), we have a canonical homomorphism
\[
  \Gamma(\rho):\Gamma(X',\mathcal{F}'(n))
  \longrightarrow\Gamma(X,\mathcal{F}(n))
\]
which gives a canonical di-homomorphism of graded modules
\[
  \Gamma_\bullet(\mathcal{F}')\longrightarrow
  \Gamma_\bullet(\mathcal{F}).
\]

Suppose that the ideal $S_+$ is generated by $S_1$, and that
$\mathcal{F}'=\widetilde{M'}$, so that $\mathcal{F}=\widetilde{M}$ with
$M=M'\otimes_{A'}A$. If $f'$ is homogeneous in $S'_+$ and
$f=\varphi(f')$, we have seen that
$M_{(f)}=M'_{(f')}\otimes_{A'}A$, and the diagram
\[
  \xymatrix{
    M'_0\ar[r]\ar[d]
    & M'_{(f')}\ar[d]
    & =\Gamma(D_+(f'),\widetilde{M'})\\
    M_0\ar[r]
    & M_{(f)}
    & =\Gamma(D_+(f),\widetilde{M})
  }
\]
is therefore commutative; from this remark and from the definition of the
homomorphism $\alpha$ (\hyperref[II.2.6.2]{2.6.2}) we immediately conclude
that the diagram
\[
\label{II.2.8.13.1}
  \xymatrix{
    M'\ar[r]^{\alpha_{M'}}\ar[d]
    & \Gamma_\bullet(\widetilde{M'})\ar[d]\\
    M\ar[r]_{\alpha_M}
    & \Gamma_\bullet(\widetilde{M})
  }
\tag{2.8.13.1}
\]
is commutative. Likewise, the diagram
\[
\label{II.2.8.13.2}
  \xymatrix{
    \widetilde{\Gamma_\bullet(\mathcal{F}')}\ar[r]^{\beta_{\mathcal{F}'}}
      \ar[d]
    & \mathcal{F}'\ar[d]\\
    \widetilde{\Gamma_\bullet(\mathcal{F})}\ar[r]_{\beta_{\mathcal{F}}}
    & \mathcal{F}
  }
\tag{2.8.13.2}
\]
is commutative (the right-hand vertical arrow being the canonical
$\Phi$-morphism
$\mathcal{F}'\to\Phi^*(\mathcal{F}')=\mathcal{F}$).
\end{env}

\oldpage[II]{47}

\begin{env}[2.8.14]
\label{II.2.8.14}
Still keeping the hypotheses and notation of
(\hyperref[II.2.8.9]{2.8.9}), let $N'$ be a second graded $S'$-module, and
let $N=N'\otimes_{A'}A$. It is immediate that the canonical di-homomorphisms
$M'\to M$, $N'\to N$ give a di-homomorphism
$M'\otimes_{S'}N'\to M\otimes_S N$ (relative to the canonical ring
homomorphism $S'\to S$), and consequently also an $S$-homomorphism of
degree~$0$
\[
  (M'\otimes_{S'}N')\otimes_{A'}A\longrightarrow M\otimes_S N,
\]
to which there corresponds (taking into account
(\hyperref[II.2.8.10]{2.8.10})) a homomorphism of
$\mathcal{O}_X$-modules
\[
\label{II.2.8.14.1}
  \Phi^*(\widetilde{M'\otimes_{S'}N'})
  \longrightarrow\widetilde{M\otimes_S N}.
\tag{2.8.14.1}
\]

Moreover, we verify immediately that the diagram
\[
\label{II.2.8.14.2}
  \xymatrix{
    \Phi^*(\widetilde{M'}\otimes_{\mathcal{O}_{X'}}\widetilde{N'})
      \ar[r]^{\sim}\ar[d]_{\Phi^*(\lambda)}
    & \widetilde{M}\otimes_{\mathcal{O}_X}\widetilde{N}
      \ar[d]^{\lambda}
    & =\Phi^*(\widetilde{M'})\otimes_{\mathcal{O}_X}
      \Phi^*(\widetilde{N'})\\
    \Phi^*(\widetilde{M'\otimes_{S'}N'})\ar[r]
    & \widetilde{M\otimes_S N}
  }
\tag{2.8.14.2}
\]
is commutative, the first row being the canonical isomorphism
(\hyperref[0.4.3.3]{0, 4.3.3}). If the ideal $S'_+$ is generated by
$S'_1$, it is clear that $S_+$ is generated by $S_1$, and the two vertical
arrows of (\hyperref[II.2.8.14.2]{2.8.14.2}) are then isomorphisms
(\hyperref[II.2.5.13]{2.5.13}); the same is therefore true of
(\hyperref[II.2.8.14.1]{2.8.14.1}).

Likewise, we have a canonical di-homomorphism
$\operatorname{Hom}_{S'}(M',N')\to\operatorname{Hom}_S(M,N)$, by assigning to
a homomorphism $u'$ of degree $k$ the homomorphism $u'\otimes1$, which is
also of degree $k$; from this we again deduce an $S$-homomorphism of
degree~$0$
\[
  (\operatorname{Hom}_{S'}(M',N'))\otimes_{A'}A
  \longrightarrow\operatorname{Hom}_S(M,N)
\]
whence a homomorphism of $\mathcal{O}_X$-modules
\[
\label{II.2.8.14.3}
  \Phi^*(\widetilde{\operatorname{Hom}_{S'}(M',N')})
  \longrightarrow\widetilde{\operatorname{Hom}_S(M,N)}.
\tag{2.8.14.3}
\]

Moreover, the diagram
\[
  \xymatrix{
    \Phi^*(\widetilde{\operatorname{Hom}_{S'}(M',N')})\ar[r]
      \ar[d]_{\Phi^*(\mu)}
    & \widetilde{\operatorname{Hom}_S(M,N)}\ar[d]^{\mu}\\
    \Phi^*(\mathcal{H}om_{\mathcal{O}_{X'}}
      (\widetilde{M'},\widetilde{N'}))\ar[r]
    & \mathcal{H}om_{\mathcal{O}_X}(\widetilde{M},\widetilde{N})
  }
\]
is commutative (the second horizontal row being the canonical homomorphism
(\hyperref[0.4.4.6]{0, 4.4.6})).
\end{env}

\oldpage[II]{48}

\begin{env}[2.8.15]
\label{II.2.8.15}
The notation and hypotheses being those of
(\hyperref[II.2.8.1]{2.8.1}), it follows from
(\hyperref[II.2.4.7]{2.4.7}) that the morphism $\Phi$ does not change, up
to isomorphisms, if we replace $S_0$ and $S'_0$ by $\mathbf{Z}$ and
$\varphi_0$ by the identity map, and then replace $S$ and $S'$ by $S^{(d)}$
and ${S'}^{(d)}$ respectively ($d>0$) and $\varphi$ by its restriction
$\varphi^{(d)}$ to $S^{(d)}$.
\end{env}

\subsection{Closed subschemes of a scheme $\operatorname{Proj}(S)$.}
\label{subsection:II.2.9}

\begin{env}[2.9.1]
\label{II.2.9.1}
If $\varphi:S\to S'$ is a homomorphism of graded rings, we say that
$\varphi$ is (TN)-\emph{surjective} (resp. (TN)-\emph{injective},
(TN)-\emph{bijective}) if there exists an integer $n$ such that, for
$k\geq n$, $\varphi_k:S_k\to S'_k$ is surjective (resp. injective,
bijective). It then follows from (\hyperref[II.2.8.15]{2.8.15}) that the
study of $\Phi$ can be reduced to the case where $\varphi$ is
\emph{surjective} (resp. \emph{injective}, \emph{bijective}). Instead of
saying that $\varphi$ is (TN)-bijective, we also say that it is then a
(TN)-\emph{isomorphism}.
\end{env}

\begin{proposition}[2.9.2]
\label{II.2.9.2}
Let $S$ be a ring graded in positive degrees and let
$X=\operatorname{Proj}(S)$.
\begin{enumerate}
  \item[\rm(i)] If $\varphi:S\to S'$ is a (TN)-surjective homomorphism of
    graded rings, the corresponding morphism $\Phi$
    (\hyperref[II.2.8.1]{2.8.1}) is defined on the whole of
    $\operatorname{Proj}(S')$ and is a closed immersion of
    $\operatorname{Proj}(S')$ into $X$. If $\mathfrak{J}$ is the kernel of
    $\varphi$, the closed subscheme of $X$ associated to $\Phi$ is defined by
    the quasi-coherent sheaf of ideals $\widetilde{\mathfrak{J}}$ of
    $\mathcal{O}_X$.
  \item[\rm(ii)] Suppose moreover that the ideal $S_+$ is generated by
    finitely many homogeneous elements of degree~$1$. Let $X'$ be a closed
    subscheme of $X$, defined by a quasi-coherent sheaf of ideals
    $\mathcal{J}$ of $\mathcal{O}_X$. Let $\mathfrak{J}$ be the graded ideal
    of $S$ that is the inverse image of $\Gamma_\bullet(\mathcal{J})$ under the
    canonical homomorphism $\alpha:S\to\Gamma_\bullet(\mathcal{O}_X)$
    (\hyperref[II.2.6.2]{2.6.2}), and set
    $S'=S/\mathfrak{J}$. Then $X'$ is the subscheme associated to the closed
    immersion $\operatorname{Proj}(S')\to X$ corresponding to the canonical
    homomorphism of graded rings $S\to S'$.
\end{enumerate}
\end{proposition}

\begin{proof}
\begin{enumerate}
  \item[\rm(i)] We may suppose $\varphi$ surjective
    (\hyperref[II.2.9.1]{2.9.1}). Since by hypothesis
    $\varphi(S_+)$ generates $S'_+$, we have
    $G(\varphi)=\operatorname{Proj}(S')$. On the other hand, the second
    assertion can be verified locally on $X$; so let $f$ be a homogeneous
    element of $S_+$ and set $f'=\varphi(f)$. Since $\varphi$ is a surjective
    homomorphism of graded rings, we verify immediately that
    $\varphi_{(f')}:S_{(f)}\to S'_{(f')}$ is surjective and that its kernel is
    $\mathfrak{J}_{(f)}$, which completes the proof of (i)
    (\hyperref[I.4.2.3]{I, 4.2.3}).
  \item[\rm(ii)] By virtue of (i), we are reduced to verifying that the
    homomorphism $\widetilde{j}:\widetilde{\mathfrak{J}}\to\mathcal{O}_X$
    deduced from the canonical injection $j:\mathfrak{J}\to S$ is an
    isomorphism of $\widetilde{\mathfrak{J}}$ onto $\mathcal{J}$, which
    follows from (\hyperref[II.2.7.11]{2.7.11}).
\end{enumerate}
\end{proof}

We note that $\mathfrak{J}$ is the \emph{largest} of the graded ideals
$\mathfrak{J}'$ of $S$ such that
$\widetilde{j}(\widetilde{\mathfrak{J}'})=\mathcal{J}$, since we verify
immediately, by going back to the definitions
(\hyperref[II.2.6.2]{2.6.2}), that this relation implies
$\alpha(\mathfrak{J}')\subset\Gamma_\bullet(\mathcal{J})$.

\begin{corollary}[2.9.3]
\label{II.2.9.3}
Suppose that the hypotheses of
(\hyperref[II.2.9.2]{2.9.2, (i)}) are satisfied and moreover that the ideal
$S_+$ is generated by $S_1$; then $\Phi^*(\widetilde{S(n)})$ is canonically
isomorphic to $\widetilde{S'(n)}$ for every $n\in\mathbf{Z}$, and consequently
$\Phi^*(\mathcal{F}(n))$ to $\Phi^*(\mathcal{F})(n)$ for every
$\mathcal{O}_X$-module $\mathcal{F}$.
\end{corollary}

\begin{proof}
This is a particular case of (\hyperref[II.2.8.8]{2.8.8}), taking into
account (\hyperref[II.2.5.10.2]{2.5.10.2}).
\end{proof}

\begin{corollary}[2.9.4]
\label{II.2.9.4}
Suppose that the hypotheses of
(\hyperref[II.2.9.2]{2.9.2, (ii)}) are satisfied. For the closed
sub-prescheme $X'$ of $X$ to be integral, it is necessary and sufficient that
the graded ideal $\mathfrak{J}$ be prime in $S$.
\end{corollary}

\oldpage[II]{49}

\begin{proof}
Since $X'$ is isomorphic to $\operatorname{Proj}(S/\mathfrak{J})$, the
condition is sufficient by virtue of (\hyperref[II.2.4.4]{2.4.4}). To see
that it is necessary, consider the exact sequence
$0\to\mathcal{J}\to\mathcal{O}_X\to\mathcal{O}_X/\mathcal{J}$, which gives an
exact sequence
\[
  0\longrightarrow\Gamma_\bullet(\mathcal{J})
  \longrightarrow\Gamma_\bullet(\mathcal{O}_X)
  \longrightarrow\Gamma_\bullet(\mathcal{O}_X/\mathcal{J}).
\]
It suffices to prove that if $f\in S_m$, $g\in S_n$ are such that the image of
$\alpha_{n+m}(fg)$ in $\Gamma_\bullet(\mathcal{O}_X/\mathcal{J})$ is zero,
then one of the images of $\alpha_m(f)$, $\alpha_n(g)$ is zero. Now, by
definition, these images are sections of the invertible
$(\mathcal{O}_X/\mathcal{J})$-modules
$\mathcal{L}=(\mathcal{O}_X/\mathcal{J})(m)$ and
$\mathcal{L}'=(\mathcal{O}_X/\mathcal{J})(n)$ over the integral scheme
$X'$; the hypothesis implies that the product of these two sections is zero
in $\mathcal{L}\otimes\mathcal{L}'$
((\hyperref[II.2.9.3]{2.9.3}) and
(\hyperref[II.2.5.14.1]{2.5.14.1})), so one of them is zero by virtue of
(\hyperref[I.7.4.4]{I, 7.4.4}).
\end{proof}

\begin{corollary}[2.9.5]
\label{II.2.9.5}
Let $A$ be a ring, $M$ an $A$-module, $S$ a graded $A$-algebra generated by
the set $S_1$ of homogeneous elements of degree~$1$, $u:M\to S_1$ a
surjective homomorphism of $A$-modules, and
$\overline{u}:\mathbf{S}(M)\to S$ the homomorphism (of $A$-algebras) from the
symmetric algebra $\mathbf{S}(M)$ of $M$ to $S$ that extends $u$. The
morphism corresponding to $\overline{u}$ is then a closed immersion of
$\operatorname{Proj}(S)$ into $\operatorname{Proj}(\mathbf{S}(M))$.
\end{corollary}

\begin{proof}
Indeed, $\overline{u}$ is surjective by hypothesis and it suffices to apply
(\hyperref[II.2.9.2]{2.9.2}).
\end{proof}

